\documentclass[11pt,a4paper]{article}
\begin{document}

\newpage
\section*{\centering Module 2: 현대 은행업 (Modern Banking)}
\addcontentsline{toc}{section}{Module 2: 현대 은행업 (Modern Banking)}

\begin{summarybox}{개요: Module 2 핵심 요약}
이 모듈은 현대 은행의 비즈니스 모델을 심층적으로 탐구합니다. 미국 은행업의 간략한 역사부터 시작하여, 현대 은행의 조직과 기능(상업 은행 및 투자 은행)을 상세히 설명합니다. 마지막으로, 현대 은행의 재무제표 분석을 통해 이들 금융기관의 운영 방식을 이해하는 것을 목표로 합니다. 이 모듈을 마치면 현대 은행의 기원, 역할, 그리고 수익 창출 방식을 명확히 이해할 수 있을 것입니다.
\end{summarybox}

\newpage
\section*{용어 정리}
\addcontentsline{toc}{section}{용어 정리}

\begin{adjustbox}{width=\textwidth,center}
\begin{tabular}{lllc}
\toprule
\textbf{용어 (Term)} & \textbf{쉬운 설명 (Easy Explanation)} & \textbf{원어 (Original)} & \textbf{비고 (Notes)} \\
\midrule
텔러 & 은행 창구에서 고객을 응대하는 직원 & Teller & 과거에는 흔했지만 지금은 드물다. \\
유닛 뱅킹 & 은행이 단 하나의 물리적 지점만 가질 수 있는 형태 & Unit Banking & 일리노이주에서 1985년까지 시행. \\
지점 네트워크 & 은행이 여러 물리적 지점을 운영하는 형태 & Branch Network & 캘리포니아주에서 일찍이 허용. \\
은행 공황 & 예금자들이 은행 파산을 우려하여 예금을 대량 인출하는 현상 & Banking Panic & 예금 보험이 없던 시절 흔했다. \\
글래스-스티걸 법 & 1933년 제정된 법으로, 상업 은행과 투자 은행 업무를 분리 & Glass-Steagall Act & 은행 시스템 안정화에 기여했으나 비효율성 비판도 받음. \\
연방 예금 보험 공사 & 은행 파산 시 예금자에게 예금을 보장해주는 기관 & Federal Deposit Insurance Corporation (FDIC) & 예금자 보호를 통해 은행 공황 방지. \\
규제 Q & 예금 계좌에 대한 이자 지급을 금지했던 규제 & Regulation Q & 은행 간 과도한 경쟁을 막기 위함. \\
은행 지주 회사 & 다른 은행이나 비은행 자회사를 소유하는 껍데기 회사 & Bank Holding Company (BHC) & 지점 제한 규제를 우회하는 데 사용. \\
규제 차익 거래 & 법의 허점을 이용하여 규제를 회피하는 행위 & Regulatory Arbitrage & BHC가 지점 제한을 우회한 사례. \\
3-6-3 규칙 & 예금 3\% 이자, 대출 6\% 이자, 오후 3시 골프를 치는 은행업 모델 & 3-6-3 Rule & 규제 보호 하에 은행이 누리던 독점적 이익. \\
머니 마켓 뮤추얼 펀드 & 은행 예금과 유사하지만 규제 Q의 적용을 받지 않는 투자 상품 & Money Market Mutual Fund & 은행 예금의 대체재로 등장. \\
증권화 & 은행 대출을 묶어 유가증권으로 만들어 판매하는 과정 & Securitization & 은행 대출의 새로운 자금 조달 모델. \\
리가-닐 법 & 은행의 주(州) 경계를 넘는 지리적 확장을 허용한 법 & Riegle-Neal Act & 은행 통합 및 효율성 증대. \\
그램-리치-블라일리 법 & 글래스-스티걸 법을 폐지하여 상업/투자 은행 업무 통합을 허용한 법 & Gramm-Leach-Bliley Act & 금융 대형화 및 복합 금융 서비스 등장. \\
금융 지주 회사 & 은행, 증권, 보험 등 다양한 금융 서비스를 제공하는 대형 금융 그룹 & Financial Holding Company & 현대 메가 은행의 조직 형태. \\
소매 금융 & 개인 및 소규모 기업 고객을 대상으로 하는 은행 서비스 & Retail Banking & 예금, 주택 담보 대출, 신용 카드 등. \\
도매 금융 & 대기업 및 기관 투자자를 대상으로 하는 은행 서비스 & Wholesale Banking & 기업 대출, 투자 은행 서비스 등. \\
커뮤니티 은행 & 지역 사회의 개인 및 소규모 기업을 위한 소규모 은행 & Community Bank & 자산 10억 달러 미만. \\
지역/초지역 은행 & 커뮤니티 은행보다 크고 전국적인 지점망을 가진 은행 & Regional/Super-regional Bank & 자산 수백억~수조 달러. \\
머니 센터 은행 & 도매 자금 조달에 크게 의존하는 대형 은행 & Money Center Bank & 뉴욕 등 주요 금융 중심지에 위치. \\
거래 계좌 & 수시 입출금이 가능하고 이자가 없는 예금 계좌 & Transaction Account & 요구불 예금 (Checking Account). \\
비거래 계좌 & 저축성 예금으로, 이자가 지급되지만 인출 제한이 있는 계좌 & Non-transaction Account & 저축 예금, 소액 양도성 예금 증서 (CD). \\
연방 기금 시장 & 은행 간 단기 무담보 대출이 이루어지는 시장 & Federal Funds Market & 은행 간 자금 부족 해소. \\
환매 조건부 채권 & 증권을 담보로 단기 자금을 빌리는 계약 & Repurchase Agreement (Repo) & 담보 대출의 일종. \\
은행 자본 & 은행 소유주의 투자금으로, 주식 및 유보 이익으로 구성 & Bank Capital & 손실 흡수 능력 제공, 가장 안정적인 자금원. \\
소비자 신용 & 개인 고객에게 제공되는 대출 & Consumer Credit & 주택 담보 대출, 자동차 대출, 학자금 대출 등. \\
상업 및 산업 대출 & 기업 고객에게 제공되는 대출 & Commercial and Industrial (C\&I) Lending) & 상업용 부동산 대출, 기간 대출, 사업 신용 한도 등. \\
관계형 대출 & 은행이 고객과의 장기적인 관계를 바탕으로 제공하는 대출 & Relationship Lending & 중견 기업 대출에 주로 적용. \\
대출 개시 & 대출 신청을 심사하고 대출 조건을 협상하여 승인하는 과정 & Origination & 대출 과정의 첫 단계. \\
대출 자금 조달 & 승인된 대출에 필요한 자금을 실제로 제공하는 과정 & Financing & 은행의 자금원 활용. \\
대출 관리 & 대출 실행 후 상환금을 수취하고 차입자의 상태를 모니터링하는 과정 & Servicing & 대출의 지속적인 관리. \\
개시 및 보유 모델 & 은행이 대출을 개시하고 자금을 조달하며 만기까지 보유하는 전통적 모델 & Originate and Hold Model & 모든 과정을 은행이 직접 수행. \\
개시 및 분배 모델 & 은행이 대출을 개시한 후 다른 기관에 판매하거나 증권화하는 모델 & Originate and Distribute Model & 수수료 수입에 중점, 효율성 증대. \\
핀테크 & 금융 기술을 의미하며, 새로운 기술을 활용한 금융 서비스 & Fintech & 온라인 대출 등으로 은행업을 혁신. \\
\bottomrule
\end{tabular}
\end{adjustbox}

\newpage
\section{Lesson 2-0: 개요 (Overview)}
\addcontentsline{toc}{section}{Lesson 2-0: 개요 (Overview)}

안녕하세요, 금융기관 모듈 2에 오신 것을 환영합니다. 이 모듈은 \textbf{현대 은행업(Modern Banking)}이라고 불립니다.

\begin{summarybox}{모듈 2의 학습 목표}
이 모듈에서는 현대 은행의 비즈니스 모델을 심층적으로 다룰 것입니다.
\begin{itemize}
    \item \textbf{미국 은행업의 간략한 역사:} 현대 은행업을 완전히 이해하기 위해 그 기원을 살펴봅니다.
    \item \textbf{현대 은행의 조직과 기능:} 상업 은행과 투자 은행 사업 모델을 분리하여 각각을 분석합니다.
    \item \textbf{재무제표 분석:} 은행의 대차대조표(Balance Sheet)와 손익계산서(Income Statement)를 면밀히 검토하여 이들 금융기관에서 실제로 어떤 일이 벌어지고 있는지 이해합니다.
\end{itemize}
이 모듈을 마치면 현대 은행이 어디에서 왔고, 무엇을 하며, 어떻게 이익을 창출하는지 명확하게 이해할 수 있을 것입니다.
\end{summarybox}

\newpage
\section{Lesson 2-1: 미국 은행업의 간략한 역사 (A Brief History of Banking in the United States)}
\addcontentsline{toc}{section}{Lesson 2-1: 미국 은행업의 간략한 역사 (A Brief History of Banking in the United States)}

\subsection{Lesson 2-1.1: 미국 초기 은행업 (Early Banking in the United States)}
\addcontentsline{toc}{subsection}{Lesson 2-1.1: 미국 초기 은행업 (Early Banking in the United States)}

우리가 '은행 시스템'이라고 부르는 것은 지난 100년 동안 극적으로 변화했습니다.

\begin{itemize}
    \item \textbf{과거:} 우리 부모님 세대는 아마도 동네 은행을 방문하여 텔러(Teller)와 직접 대화했을 것입니다. 텔러는 은행에서 일하는 실제 사람을 의미합니다.
    \item \textbf{현재:} 요즘에는 ATM조차 거의 사용하지 않습니다. 대부분의 은행 업무는 온라인이나 스마트폰을 통해 이루어집니다.
\end{itemize}

은행업의 기술적, 경쟁적 환경이 변화함에 따라 은행 산업의 구조도 변화했습니다. 여기서 '구조'는 두 가지를 의미합니다.
\begin{enumerate}
    \item \textbf{은행의 비즈니스 모델:} 은행이 판매하는 상품의 종류를 포함합니다.
    \item \textbf{은행의 수, 규모, 물리적 위치:} 은행의 지리적 분포를 의미합니다.
\end{enumerate}
이번 세션에서는 미국 은행 산업의 역사적 추세에 초점을 맞출 것입니다. 기술과 규제의 변화가 21세기 은행의 모습에 어떻게 영향을 미쳤는지 살펴보겠습니다.

\subsubsection*{전통적인 은행의 비즈니스 모델}
전통적인 은행은 여러 가지 일을 하지만, 그 비즈니스 모델은 크게 다음으로 요약할 수 있습니다.
\begin{itemize}
    \item \textbf{예금 수취:} 예금자로부터 돈(달러)을 받아 은행 계좌를 제공합니다.
    \item \textbf{자금 대출:} 이 자금을 주로 대출 형태로 차입자에게 전달합니다.
\end{itemize}

\begin{examplebox}{은행 계좌의 가치}
은행 계좌는 기본적으로 돈과 같습니다. 우리는 은행에 우리의 자산을 안전하게 보관할 수 있으며, 가격 변동 위험 없이 결제 시스템에 접근할 수 있습니다. 예를 들어, 수표를 발행하거나 직불 카드(Debit Card)를 사용할 수 있습니다.
\end{examplebox}

은행은 또한 신용도 높은 차입자에게만 대출이 이루어지도록 함으로써 가치를 더합니다. 전통적인 은행은 예금을 빌려주고, 대출 이자 수입과 예금 이자 지급의 차이(스프레드)를 통해 수익을 얻습니다. 이들은 노력과 위험 감수에 대한 보상을 받습니다.

\subsubsection*{현대 은행업의 복잡성}
현대 은행업은 매우 복잡합니다. 예를 들어, JP모건 체이스(JPMorgan Chase)의 조직도를 보면 이 회사가 운영하는 다양한 사업 부문을 알 수 있습니다. 여기에는 일반 개인 고객을 위한 \textbf{소비자 금융(Consumer Banking)} 사업뿐만 아니라 훨씬 더 많은 상품을 포함하는 \textbf{도매 금융(Wholesale Banking)} 사업도 포함됩니다. JP모건과 같은 조직은 수만 개의 자회사를 통해 모든 종류의 금융 서비스를 제공합니다.

\begin{center}
    \textit{어떻게 전통적인 은행 모델에서 JP모건 체이스와 같은 '금융 슈퍼마켓'으로 발전하게 되었을까요?}
\end{center}

\subsubsection*{은행의 조직 형태와 주요 논점}
이번 세션에서는 세 가지 주요 조직 형태를 논의할 것입니다.
\begin{enumerate}
    \item \textbf{유닛 뱅크 (Unit Bank):} 단일 지점 은행
    \item \textbf{지점 네트워크 (Branch Network):} 여러 지점을 가진 은행
    \item \textbf{은행 지주 회사 (Bank Holding Company):} 여러 은행 및 비은행 자회사를 소유하는 회사
\end{enumerate}
이 논의를 관통하는 두 가지 핵심 주제가 있습니다.
\begin{enumerate}
    \item \textbf{은행이 법적으로 운영할 수 있는 지역:} 은행이 지리적으로 어디까지 확장할 수 있는가?
    \item \textbf{은행이 판매할 수 있는 금융 상품 및 서비스의 종류:} 은행이 어떤 종류의 사업을 할 수 있는가?
\end{enumerate}
21세기에는 은행이 원하는 곳 어디에서든 운영하고 상상할 수 있는 모든 금융 서비스를 제공할 수 있는 것처럼 보입니다. 그러나 이것이 항상 그랬던 것은 아니라는 점을 인식하는 것이 중요합니다. 역사적으로 은행의 확장이나 사업 변경 허용은 격렬한 정책 논쟁의 원천이었습니다.

\subsubsection*{지리적 확장 제한: 유닛 뱅킹 vs. 지점 네트워크}
은행이 지리적으로 어떻게 분포되었는지 살펴보겠습니다. 이는 은행의 물리적인 지점 수와 위치를 의미합니다. 과거에는 은행이 가질 수 있는 지점 수에 많은 제한이 있었습니다. 주(州)마다 접근 방식이 달랐습니다.

\begin{itemize}
    \item \textbf{엄격한 주 (예: 일리노이):} 일리노이주는 \textbf{유닛 뱅킹(Unit Banking)}만을 허용했습니다. 이는 은행이 단 하나의 물리적 위치만 가질 수 있다는 것을 의미했습니다.
    \begin{examplebox}{콘티넨탈 일리노이 은행 (Continental Illinois Bank)}
    시카고 금융 지구의 라살 스트리트 208번지에 위치한 콘티넨탈 일리노이 은행이 그 예입니다. CEO를 포함한 모든 직원이 한 곳에 있었고, 모든 대출 기록도 이 건물에 보관되었습니다. 고객들은 어디로 가야 할지 명확히 알았습니다. 유닛 뱅킹은 일리노이주에서 1870년부터 1985년까지 법으로 시행되었습니다.
    \end{examplebox}
    \item \textbf{자유로운 주 (예: 캘리포니아):} 캘리포니아주는 항상 은행이 주 경계 내에서 \textbf{지점 네트워크(Branch Network)}를 가질 수 있도록 허용했습니다.
    \begin{examplebox}{뱅크 오브 아메리카 (Bank of America) 지점 네트워크 (1929년경)}
    1929년경 뱅크 오브 아메리카의 지점 네트워크를 보면, 샌프란시스코에 본사를 두고 주 전역에 지점을 운영했습니다.
    \end{examplebox}
\end{itemize}

일리노이주와 다른 주들이 엄격한 지점 규제를 시행한 데에는 여러 가지 이유가 있었습니다.
\begin{itemize}
    \item \textbf{지점 확장에 반대하는 주장:} 은행이 너무 커져 독점력을 행사할 수 있다는 우려가 있었습니다. 특히 소도시의 은행가들은 대도시 은행과의 경쟁으로부터 자신들을 보호하기 위해 지점 제한을 강력히 지지했습니다.
    \item \textbf{지점 확장에 찬성하는 주장:} 다른 주들은 더 크고 지리적으로 분산된 은행이 더 회복력 있고 안전할 것이라고 믿었습니다.
\end{itemize}
유닛 뱅킹과 지점 네트워크 중 어느 것이 더 안전한지에 대한 의견 불일치가 많았지만, 이 시기에는 \textbf{은행 공황(Banking Panics)}이 흔했습니다.

\begin{examplebox}{은행 공황에 대한 유명한 인용문}
\begin{itemize}
    \item "그때는 은행 공황이 얼마나 자주 일어났는지"를 지적하는 첫 번째 인용문.
    \item "결국 전 세계적인 범위로 확대되었다"는 두 번째 인용문.
\end{itemize}
\end{examplebox}

공황은 다음과 같이 작동합니다. 일반적으로 일부 예금자들이 저축을 잃을까 봐 걱정하기 시작합니다. 왜 이런 일이 일어났을까요? 예금자들은 경제가 악화될 수 있고, 자신들의 은행이 대출에서 큰 손실을 입어 결국 파산할 수 있다고 우려하기 시작했습니다. 이때는 \textbf{예금 보험(Deposit Insurance)}이 없던 시절이었습니다. 이는 제때 돈을 인출하지 못한 예금자들만 손실을 입게 된다는 것을 의미했습니다. 따라서 '뱅크 런(Run on the Bank)'이 발생했습니다.

이러한 은행 공황은 1930년대 \textbf{대공황(Great Depression)} 동안 절정에 달했습니다. 이는 아마도 현대 역사상 최악의 글로벌 금융 위기였을 것입니다. 미국 내 모든 은행의 3분의 1 이상이 파산했고, 수백만 가구가 평생 모은 저축을 잃었으며, 경제는 붕괴했습니다.

\subsubsection*{1933년 은행법 (글래스-스티걸 법)}
은행 위기와 그에 따른 대공황에 대한 비상 대응으로, 프랭클린 루즈벨트 대통령은 \textbf{1933년 은행법(Banking Act of 1933)}을 제정했습니다. 이 유명한 법안은 \textbf{글래스-스티걸 법(Glass-Steagall Act)}으로도 알려져 있습니다.

\begin{summarybox}{글래스-스티걸 법의 주요 변화}
\begin{itemize}
    \item \textbf{주(州) 간 은행 지점 설치 금지 (Ban on Interstate Bank Branching):} 은행이 주 경계를 넘어 자회사를 가질 수 없도록 하는 지리적 제한이었습니다.
    \item \textbf{투자 은행 업무 금지 (Ban on Investment Banking):} 상업 은행은 사모 증권의 인수(Underwriting)나 거래, 심지어 보험 판매도 금지되었습니다.
    \item \textbf{규제 Q (Regulation Q) 포함:} 예금 계좌에 대한 이자 지급을 금지했습니다.
    \item \textbf{연방 예금 보험 공사 (Federal Deposit Insurance Corporation, FDIC) 설립:} 은행이 파산할 경우 개별 예금자에게 예금을 보장하여 저축을 잃지 않도록 했습니다.
\end{itemize}
\end{summarybox}

이러한 금지 조치들은 은행 공황의 원인에 대한 당시의 믿음을 반영했습니다.
\begin{itemize}
    \item 주(州) 간 지점 확장은 공황을 들불처럼 퍼뜨린다고 보았습니다.
    \item 투자 은행 업무는 너무 위험하다고 여겨졌습니다.
    \item 예금에 대한 이자 지급은 은행 간 과도한 경쟁을 유발한다고 생각했습니다.
    \item 뱅크 런은 저축을 잃을까 봐 걱정하는 혼란스러운 투자자들 때문에 발생한다고 보았습니다.
\end{itemize}
글래스-스티걸 법은 고도로 규제되고 보호받는 상업 은행 산업을 만들었습니다.

\begin{figure}[h!]
    \centering
    \includegraphics[width=0.8\textwidth]{example_graph_banks_1935_1956.png} % 가상의 이미지 파일명
    \caption{1935년에서 1956년 사이 미국 내 보험 가입 상업 은행 수 (단위 은행 및 다중 지점 은행)}
    \label{fig:banks_1935_1956}
    \textit{이 그래프는 1935년에서 1956년 사이의 보험 가입 상업 은행 수를 보여줍니다. 이 기간 동안 은행의 수는 많고 안정적이었으며, 대부분은 지역 독점권을 가진 유닛 뱅크였습니다. 이는 정부가 보장하는 은행 시스템으로, 단순한 상품 제공과 제한된 위험 감수 기회를 가졌습니다.}
\end{figure}

글래스-스티걸 법은 은행 시스템에 대한 신뢰를 회복시켰습니다.

\begin{figure}[h!]
    \centering
    \includegraphics[width=0.8\textwidth]{example_graph_bank_failures.png} % 가상의 이미지 파일명
    \caption{미국 내 은행 파산 건수 (1933년 이후)}
    \label{fig:bank_failures}
    \textit{이 차트는 은행 파산 건수가 1933년 약 4,000건에서 효과적으로 0건으로 급감했음을 보여줍니다. 그러나 당시 많은 비평가들은 고도로 규제된 미국 은행 시스템이 너무 파편화되어 있고 비효율적이라고 보았습니다.}
\end{figure}

\newpage
\subsection{Lesson 2-1.2: 은행 지주 회사 (Bank Holding Companies)}
\addcontentsline{toc}{subsection}{Lesson 2-1.2: 은행 지주 회사 (Bank Holding Companies)}

다음으로 등장한 지배적인 조직 형태는 \textbf{은행 지주 회사(Bank Holding Company, BHC)}였습니다.

\begin{itemize}
    \item \textbf{지주 회사(Holding Company):} 다른 기업이나 자회사 그룹을 소유하는 껍데기 회사(Shell Corporation)입니다.
    \item \textbf{은행 지주 회사(BHC):} 하나 이상의 은행 자회사를 소유합니다.
\end{itemize}

\begin{examplebox}{규제 회피를 위한 은행 지주 회사 활용}
1950년이라고 가정해 봅시다. 제가 샴페인(Champaign)에 매우 성공적인 은행을 소유하고 있고, 지리적으로 다각화하기 위해 세인트 클레어 카운티(Saint Clair County)로 사업을 확장하고 싶습니다. 그러나 일리노이주는 유닛 뱅킹만 허용하므로 이는 불가능합니다.

지점 제한을 회피하는 한 가지 방법은 다음과 같습니다.
\begin{enumerate}
    \item 세인트 클레어에 완전히 별개의 은행을 설립합니다.
    \item 두 은행을 모두 인수하는 지주 회사를 설립합니다.
\end{enumerate}
이러한 은행 자회사들은 독립적인 은행으로 간주될 수 있으므로 법을 준수하는 것으로 보일 수 있습니다.

더 나아가, 제 지주 회사가 일부 자산 관리 서비스를 제공하는 비은행 자회사(Non-bank Subsidiary)를 설립할 수도 있습니다. 이 비은행 자회사는 어떤 은행과도 직접적으로 관련이 없으므로, 아마도 문제가 없을 것입니다.
\end{examplebox}

이러한 유형의 행태를 종종 \textbf{규제 차익 거래(Regulatory Arbitrage)}라고 부르는데, 이는 법의 허점을 이용한 것이기 때문입니다. 1956년과 1970년의 \textbf{은행 지주 회사법(Bank Holding Company Acts)}을 통해 의회는 이러한 허점을 막으려 했지만, 이미 '고양이는 가방 밖으로 나와버린' 상태였습니다. (즉, 이미 널리 퍼져 되돌리기 어려웠다는 의미입니다.)

\begin{figure}[h!]
    \centering
    \includegraphics[width=0.8\textwidth]{example_graph_banks_1980s.png} % 가상의 이미지 파일명
    \caption{1980년대까지의 미국 내 은행 조직 형태 변화}
    \label{fig:banks_1980s}
    \textit{이 차트는 1980년대까지의 은행 조직 형태 변화를 보여줍니다. 1980년대에 이르러 은행 지주 회사를 통한 지점 확장이 지배적인 조직 형태가 되었습니다.}
\end{figure}

1935년부터 1980년까지의 기간은 미국 상업 은행업의 전성기로 여겨지기도 합니다. 은행들은 산업으로서 독점적인 권한을 가졌습니다. 이는 상업 은행업의 유명한 \textbf{3-6-3 규칙(3-6-3 Rule)}을 낳았습니다.

\begin{examplebox}{3-6-3 규칙에 대한 인용문}
"미국 은행가들은 수십 년 동안 3-6-3 규칙에 따라 운영되었습니다. 예금자에게는 3\% 이자를 지급하고, 돈은 6\%에 대출하며, 오후 3시에는 골프를 치러 갔습니다. 연방 및 주 법률이 그들이 운영하는 엄격한 규칙을 정하고 경쟁자로부터 그들을 보호했기 때문에 그들은 그렇게 정확하게 행동할 수 있었습니다."

"그 결과, 은행가들의 권력과 명성은 그들의 금고만큼이나 안전했고, 이익은 꾸준하고 확실했습니다."
\end{examplebox}

그러나 이 모든 것은 1970년대와 80년대부터 변하기 시작했습니다. "갑자기 그 모든 것이 사라졌습니다. 은행가들은 이제 대공황 이후 가장 힘든 생존 시험에 직면해 있습니다."

\begin{figure}[h!]
    \centering
    \includegraphics[width=0.8\textwidth]{example_graph_commercial_banks_peak.png} % 가상의 이미지 파일명
    \caption{미국 내 상업 은행 수의 변화 (1985년 이후)}
    \label{fig:commercial_banks_peak}
    \textit{이 차트는 미국 내 고유 상업 은행의 수가 1985년경 정점을 찍었음을 보여줍니다. 그 후에는 개별 기관의 수가 급격히 감소했습니다. 이는 극적인 통합(Consolidation)의 시기였습니다. 1999년에는 은행 수가 약 절반으로 줄어들어 약 8,000개가 되었습니다.}
\end{figure}

\subsubsection*{은행 통합의 주요 원인}
그렇다면 이러한 통합의 주요 원인은 무엇이었을까요?

\begin{enumerate}
    \item \textbf{기술 (Technology):}
    1970년대부터 우편, 전화, 팩스, 그리고 결국 인터넷이 보편화되면서 은행업의 지역적 특성이 덜 중요해졌습니다. 우리는 더 이상 동네 은행에 직접 가서 대출 담당자와 이야기할 필요가 없어졌습니다. 요즘에는 스마트폰으로 예금, 직불 카드, 주택 담보 대출 상환 등 모든 것을 관리할 수 있습니다. 이러한 변화는 1970년대에 시작되었습니다.

    \item \textbf{채권 시장의 금융 혁신 (Financial Innovations in Debt Markets):}
    이는 예금을 유치하고 대출을 제공하는 전통적인 은행의 능력에 대한 공격이었습니다.
    \begin{itemize}
        \item \textbf{예금 대체재 등장:} 앞서 언급했듯이, 규제 Q(Regulation Q) 하에서는 은행이 예금 계좌에 이자를 지급하는 능력이 제한되었습니다. 1980년대부터 \textbf{머니 마켓 뮤추얼 펀드(Money Market Mutual Funds)}가 등장했습니다. 이 뮤추얼 펀드는 비은행 기관으로, 예금과 유사한 계좌를 제공하지만 규제 Q의 적용을 받지 않아 매우 안전하고 편리했습니다.
        \item \textbf{은행 대출 대체재 등장:} 1980년대에는 공공 채권 시장에서 엄청난 변화가 있었습니다. 기업 금융 측면에서는 단기 기업 어음(Commercial Paper)과 장기 정크 본드(Junk Bond) 시장이 등장했습니다. 그리고 \textbf{증권화(Securitization)}가 은행 대출을 위한 완전히 새로운 자금 조달 모델을 만들어냈습니다. 이는 은행 산업에 더 큰 혼란을 야기했습니다.
    \end{itemize}
\end{enumerate}

이러한 변화는 은행에 대한 주요 경쟁 압력의 원천이었습니다. 지역 은행 독점은 붕괴되었습니다. 이로 인해 은행들은 어려운 선택에 직면했습니다.
\begin{itemize}
    \item 사업을 중단하거나,
    \item 유기적으로 규모를 확장하거나,
    \item 합병 또는 인수를 통해 규모를 확장하는 것이었습니다.
\end{itemize}
규모를 확장함으로써 은행들은 \textbf{규모의 경제(Economies of Scale)}와 \textbf{범위의 경제(Economies of Scope)}를 활용하여 중복 비용을 제거할 수 있었습니다.

\newpage
\subsection{Lesson 2-1.3: 규제 완화와 현대 은행업 (Deregulation and Modern Banking)}
\addcontentsline{toc}{subsection}{Lesson 2-1.3: 규제 완화와 현대 은행업 (Deregulation and Modern Banking)}

이러한 경쟁의 시기는 은행 규제에 대한 패러다임 전환으로 이어졌습니다. 입법자들과 규제 당국은 과도한 은행 위험 감수나 독점적 은행에 대한 우려를 덜게 되었습니다. 대신, 경쟁에 실패하여 문을 닫는 은행에 대해 더 걱정하게 되었습니다.

상업 은행 산업을 현대화한 두 가지 주요 법안이 있었습니다. 둘 다 빌 클린턴 대통령에 의해 법으로 서명되었습니다. 이러한 규제 완화는 은행이 규모의 경제와 범위의 경제를 활용하여 효율성을 높이는 데 도움이 되도록 설계되었습니다.

\begin{enumerate}
    \item \textbf{리가-닐 법 (Riegle-Neal Act):}
    이 법은 은행이 주(州) 경계를 넘어 운영함으로써 지리적으로 다각화할 수 있도록 허용했습니다. 이제 은행들은 미국 전역으로 사업을 자유롭게 확장할 수 있게 되었습니다. 예를 들어, 다른 주에 있는 은행을 인수하는 방식으로 말입니다. \textbf{유닛 뱅킹(Unit Banking)}은 과거의 일이 되었습니다.

    \item \textbf{그램-리치-블라일리 법 (Gramm-Leach-Bliley Act):}
    이 법은 1933년 글래스-스티걸 법에 있던 투자 은행 업무 금지 조항을 폐지했습니다. 이제 상업 은행업, 투자 은행업, 보험 인수(Insurance Underwriting)가 모두 한 지붕 아래에서 이루어질 수 있게 되었습니다. 이는 은행이 제품 라인을 다각화할 수 있도록 허용했습니다.
\end{enumerate}

\begin{examplebox}{제품 라인 다각화의 이점: 정보 비용 절감}
은행은 정보 사업을 합니다. 여러 제품 라인에 걸쳐 운영하면 정보 비용을 낮출 수 있습니다.
\begin{itemize}
    \item 예를 들어, 고객이 당좌 예금 계좌에 꾸준한 수입 흐름을 가지고 있다는 것을 안다면, 은행은 그 고객에게 생명 보험에 대한 더 좋은 조건을 제시할 의향이 있을 수 있습니다.
    \item 또는 과거에 대출을 해준 기업 고객이라면, 다음 채권 발행을 인수하는 데 더 유리한 위치에 있을 수 있습니다.
\end{itemize}
\end{examplebox}

이러한 규제 완화는 또 다른 통합과 재편성을 촉발했습니다. 여기에는 JP모건 체이스, 뱅크 오브 아메리카, 웰스 파고와 같은 세계 최대 은행들을 탄생시킨 대규모 합병(Mega Mergers)의 물결이 포함되었습니다.

\begin{figure}[h!]
    \centering
    \includegraphics[width=0.8\textwidth]{example_graph_us_banks_final.png} % 가상의 이미지 파일명
    \caption{미국 내 은행 수 변화 (최종 시계열)}
    \label{fig:us_banks_final}
    \textit{이 차트는 미국 은행 수를 마지막으로 보여줍니다. 시계열의 끝 부분에서 지점형 은행(Branched Banks)이 이제 지배적임을 알 수 있습니다.}
\end{figure}

\begin{figure}[h!]
    \centering
    \includegraphics[width=0.8\textwidth]{example_map_bank_of_america_2010.png} % 가상의 이미지 파일명
    \caption{뱅크 오브 아메리카의 2010년 지리적 분포}
    \label{fig:bank_of_america_2010}
    \textit{마지막으로, 2010년 뱅크 오브 아메리카의 지리적 분포를 보여줍니다. 각 점은 은행 지점을 나타냅니다. 가장 큰 은행들은 이제 미국 거의 모든 곳에서 운영됩니다.}
\end{figure}

오늘날, 소수의 거대 은행들이 모든 종류의 금융 서비스를 한 지붕 아래에서 제공합니다. 이러한 기관들은 \textbf{금융 지주 회사(Financial Holding Companies)}로 조직되어 있습니다. 이들은 수천 개의 은행 및 비은행 자회사를 포함합니다.

\begin{examplebox}{JP모건 체이스의 현대적 조직 구조}
JP모건 체이스를 다시 살펴보겠습니다. 주요 활동은 다음과 같습니다.
\begin{itemize}
    \item \textbf{소매 및 상업 은행업 (Retail and Commercial Banking):} 일반 개인 및 기업 고객 대상.
    \item \textbf{투자 은행업 (Investment Banking):} 기업의 자금 조달 및 인수합병 자문.
    \item \textbf{증권 회사 기능 (Securities Firm Functions):} 브로커-딜러 활동(Broker-Dealer Activities) 포함.
    \item \textbf{자산 관리 (Asset Management):} 뮤추얼 펀드(Mutual Funds) 및 헤지 펀드(Hedge Funds) 컬렉션을 운영.
\end{itemize}
이러한 새로운 조직 구조를 통해 현대 은행들은 \textbf{수수료 수입(Fee Generation)}에 더 큰 비중을 둡니다. 투자 은행과 증권 회사들은 제공하는 서비스로부터 수수료를 얻습니다. 과거의 3-6-3 은행 모델은 이제 옛말이 되었습니다.
\end{examplebox}

\begin{summarybox}{Lesson 2-1 요약}
이번 세션에서는 20세기 초부터 현재에 이르기까지 은행업의 역사를 논의했습니다.
\begin{itemize}
    \item \textbf{초기 은행:} 단순했으며, 주로 예금 수취와 대출에 집중했고, 종종 한 곳의 물리적 위치에서 운영되었습니다. 규제는 엄격했습니다.
    \item \textbf{변화의 시작:} 시간이 지남에 따라 새로운 기술과 비은행 기관(Nonbanks)과의 경쟁이 전통적인 은행의 시장 지위를 약화시켰습니다.
    \item \textbf{규제 완화:} 은행의 현대화를 돕기 위해 규제 완화의 물결이 뒤따랐습니다. 은행들은 지리적으로 확장하고 새로운 제품 시장으로 진출하는 것이 허용되었습니다.
    \item \textbf{현대 은행의 탄생:} 이는 오늘날 우리가 보는 금융 지주 회사 구조와 집중된 은행 산업을 낳았습니다.
\end{itemize}
\textbf{질문:} 너무 많은 규제 완화가 있었을까요? 일부 거대 은행들은 '너무 커서 망할 수 없는(Too Big to Fail)' 존재가 되었을까요? 2008년 금융 위기에서 가장 큰 미국 은행들이 핵심적인 역할을 했다는 점은 그 답이 '그렇다'일 수 있음을 시사합니다. 이 복잡한 질문과 규제 당국이 이에 대해 무엇을 했는지는 나중에 다룰 것입니다.
\end{summarybox}

\newpage
\section{Lesson 2-2: 현대 은행의 조직과 기능 (Organization and Functions of Modern Banks)}
\addcontentsline{toc}{section}{Lesson 2-2: 현대 은행의 조직과 기능 (Organization and Functions of Modern Banks)}

우리는 다른 어떤 금융기관보다 은행과 더 자주 상호작용할 것입니다. 은행은 광범위한 금융 서비스를 제공합니다.
\begin{itemize}
    \item 우리의 돈을 안전하게 보관합니다.
    \item 결제 시스템에 대한 접근을 제공합니다.
    \item 금융 시장에 대한 접근을 제공합니다.
    \item 신용도 높은 차입자에게 대출을 해줍니다.
\end{itemize}
더 큰 은행들은 상업 은행 서비스와 투자 은행 서비스를 모두 한 지붕 아래에서 제공하는 '금융 슈퍼마켓'입니다. 이번 세션에서는 \textbf{상업 은행(Commercial Banks)}의 비즈니스 모델에 초점을 맞출 것입니다. 상업 은행이 어떻게 자금을 조달하고(Borrowing), 무엇을 하는지(Lending) 설명할 것입니다. 그 과정에서 \textbf{소매 금융(Retail Banking)}과 \textbf{도매 금융(Wholesale Banking)}의 차이점을 설명할 것입니다.

\subsection{Lesson 2-2.1: 은행 유형 및 자금 조달 (Bank Types and Borrowing)}
\addcontentsline{toc}{subsection}{Lesson 2-2.1: 은행 유형 및 자금 조달 (Bank Types and Borrowing)}

먼저 상업 은행이 무엇인지 정의해 봅시다.
\begin{summarybox}{상업 은행의 정의}
\textbf{상업 은행(Commercial Bank)}은 예금을 수취하고 소비자 대출, 부동산 대출, 기업 대출을 제공하는 금융기관입니다. 역사적으로 상업 은행은 기업 고객을 대상으로 했기 때문에 '상업(Commercial)'이라는 이름이 붙었지만, 요즘에는 가계와 개인에게도 서비스를 제공합니다.
\end{summarybox}

우리는 일반적으로 상업 은행을 규모와 기능에 따라 분류합니다.

\subsubsection*{상업 은행의 유형}
\begin{enumerate}
    \item \textbf{커뮤니티 은행 (Community Banks):}
    \begin{itemize}
        \item \textbf{특징:} 지역 사회의 소비자와 소규모 기업에 서비스를 제공하는 소규모 은행입니다.
        \item \textbf{운영 방식:} 개인 고객으로부터 예금을 받아 지역 사회의 개인 및 소규모 기업에 대출을 해줍니다. 이를 \textbf{소매 금융(Retail Banking)}이라고 합니다.
        \item \textbf{규모:} 주로 10억 달러 미만의 자산(대부분 대출)을 관리합니다.
    \end{itemize}

    \item \textbf{지역/초지역 은행 (Regional/Super-regional Banks):}
    \begin{itemize}
        \item \textbf{특징:} 커뮤니티 은행보다 규모가 크며, 수백억 달러에서 1조 달러 이상의 자산을 가집니다. 대도시에 본사를 두고 전국적으로 사무실과 지점을 운영할 수 있습니다.
        \item \textbf{조직 형태:} 일부는 여러 은행 및 비은행 자회사를 가진 금융 지주 회사로 조직되어 있습니다.
        \item \textbf{대출 활동:} 상업 은행 부문은 더 큰 규모로 운영됩니다. 소매 대출뿐만 아니라 대기업 고객을 대상으로 하는 상업 및 산업 대출(Commercial and Industrial Lending)과 같은 \textbf{도매 대출(Wholesale Lending)}도 합니다.
        \item \textbf{자금 조달:} 지점 네트워크를 통해 개인 예금자로부터 자금을 조달하는 것 외에도 \textbf{도매 자금 조달 시장(Wholesale Funding Markets)}에 접근할 수 있습니다. (이 빌린 자금에 대해서는 잠시 후에 더 자세히 설명하겠습니다.)
        \item \textbf{예시:} 웰스 파고(Wells Fargo)는 메가 은행이지만, 전국에서 가장 큰 지점 네트워크를 가지고 있어 예금 자금을 쉽게 확보할 수 있습니다. 따라서 웰스 파고는 초지역 은행으로 분류됩니다.
    \end{itemize}

    \item \textbf{머니 센터 은행 (Money Center Banks):}
    \begin{itemize}
        \item \textbf{특징:} 대출 활동을 지원하기 위해 \textbf{도매 자금 조달(Wholesale Funding)}에 크게 의존하는 메가 은행입니다.
        \item \textbf{위치:} 미국 기반의 머니 센터 은행은 주로 뉴욕시에 위치하며, 은행 간 대출(Interbank Loan) 및 통화 교환 시장의 중심에 있습니다.
        \item \textbf{예시:} 씨티그룹(Citigroup)은 머니 센터 은행입니다.
    \end{itemize}
\end{enumerate}

\subsubsection*{은행의 자금 조달 (Borrowing)}
상업 은행은 돈을 빌리고(Borrow), 빌려주며(Lend), 그 차이(Spread)를 벌어들입니다. 돈을 빌리고 그에 대한 비용을 지불한 다음, 돈을 빌려주고 그 투자에서 이자를 벌어들입니다. 이것이 실제로 어떤 의미일까요? 은행은 어떻게 자금을 조달하고, 어떻게 대출을 해줄까요?

먼저 \textbf{자금 조달(Funding)} 측면부터 살펴보겠습니다. 은행은 저축자로부터 자금을 얻습니다. 개인과 기업으로부터 예금을 유치합니다. 예금자들이 은행에 돈을 빌려주도록 설득하기 위해 은행은 안전성과 편리함을 제공합니다. 은행은 또한 금융 시장을 통해 기업 및 다른 금융기관으로부터 자금을 빌릴 수 있습니다. 이는 일반적으로 이자 지급을 필요로 합니다.

\begin{summarybox}{은행의 주요 자금 조달원}
\begin{enumerate}
    \item \textbf{예금 계좌 (Deposit Accounts):}
    \begin{itemize}
        \item \textbf{거래 계좌 (Transaction Accounts):}
        \begin{itemize}
            \item \textbf{요구불 예금 (Demand Deposits):} 우리의 당좌 예금 계좌(Checking Accounts)를 포함합니다. 이자는 지급되지 않지만, 원할 때 언제든지 자금을 인출할 수 있습니다.
        \end{itemize}
        \item \textbf{비거래 계좌 (Non-transaction Deposits):}
        \begin{itemize}
            \item \textbf{저축 계좌 (Savings Accounts) 및 소액 양도성 예금 증서 (Small Certificates of Deposit, CD):} 일반적으로 이자를 제공하는 저축 상품입니다. 대신 이러한 예금은 접근하기가 조금 더 어렵습니다. 한 달에 제한된 횟수만 인출할 수 있거나, CD의 경우 조기 해지 시 상당한 벌금이 부과될 수 있습니다.
        \end{itemize}
    \end{itemize}

    \item \textbf{차입 (Borrowing):}
    특히 대형 은행의 경우, 차입은 두 번째 주요 자금원입니다.
    \begin{itemize}
        \item \textbf{은행 간 시장 (Interbank Market) / 연방 기금 시장 (Federal Funds Market):} 자금 부족을 겪는 한 은행이 다른 은행으로부터 자금을 빌릴 수 있는 시장입니다. 이는 단기 무담보 대출(Short Term Unsecured Loans)입니다.
        \item \textbf{환매 조건부 채권 (Repurchase Agreement, Repo):} \textbf{담보부(Secured Basis)}로 자금을 빌릴 수 있는 금융 상품입니다. 레포는 전당포와 매우 유사하게 작동합니다.
        \begin{examplebox}{레포(Repo)의 전당포 비유}
        제가 다음 일리노이 농구 경기 티켓을 사기 위해 몇백 달러가 정말 필요하다고 가정해 봅시다. 저는 골프채를 전당포에 가져갈 수 있습니다. 우리 모두 이 골프채가 500달러 가치가 있다는 데 동의합니다. 전당포 주인은 골프채를 받고 저에게 400달러를 줄 것입니다. 이것은 골프채를 담보로 한 400달러 대출입니다. 그리고 다음 주에 월급을 받으면 500달러를 주고 골프채를 다시 살 수 있습니다. 만약 제가 다음 주에 나타나지 않으면, 전당포 주인은 제 골프채를 가지게 됩니다. 이것이 기본적으로 환매 조건부 채권(Repo)이 작동하는 방식입니다.
        \end{examplebox}
        은행이 레포를 사용할 때는 연금 기금(Pension Funds)과 같은 현금이 풍부한 기관 투자자로부터 하룻밤 동안 현금을 빌리기 위해 국채(Treasury Bills)와 같은 고품질 증권을 담보로 제공합니다.

        \textbf{레포의 문제점:} 레포와 다른 단기 차입 자금은 투자자들이 걱정하기 시작하면 사라지는 경향이 있습니다. 2007년 공황은 레포 시장의 뱅크 런이었다고 유명하게 관찰되었습니다. 미상환 레포의 규모가 거의 하룻밤 사이에 붕괴되었고, 일부 은행들은 정말 중요한 자금원을 잃었습니다. 리먼 브라더스(Lehman Brothers)가 대표적인 예입니다.
    \end{itemize}

    \item \textbf{은행 자본 (Bank Capital):}
    자금 조달 퍼즐의 마지막 조각은 은행 자본입니다. 이는 은행에 대한 소유주의 투자입니다. 주로 보통주(Common Stock)와 우선주(Preferred Stock)로 구성되며, 다른 사업과 마찬가지로 유보 이익(Retained Earnings), 즉 은행에 재투자된 이익을 포함합니다. 부채와 달리 은행 자본은 만기가 없으므로 가장 안정적인(Sticky) 자금원입니다. 그러나 주식은 가격 위험에 노출되어 있기 때문에 예금보다 조달 비용이 더 비쌉니다.
\end{enumerate}
\end{summarybox}

\begin{table}[h!]
    \centering
    \caption{은행 자금 조달원의 비용 및 안정성 비교}
    \label{tab:funding_comparison}
    \begin{adjustbox}{width=\textwidth,center}
    \begin{tabular}{p{0.25\textwidth}p{0.35\textwidth}p{0.35\textwidth}}
        \toprule
        \textbf{자금 조달원} & \textbf{비용 (Cost)} & \textbf{안정성 (Stickiness)} \\
        \midrule
        \textbf{소매 예금 (Retail Deposits)} &
        \begin{itemize}
            \item 대규모 인프라 및 고정 비용 필요
            \item 한계 비용은 저렴 (Cheap on the margin)
        \end{itemize} &
        \begin{itemize}
            \item 정부 보증 예금 보험 덕분에 공황 시에도 은행 시스템 내에 머무는 경향이 있음 (Sticky)
        \end{itemize} \\
        \addlinespace
        \textbf{도매 자금 (Wholesale Funding)} &
        \begin{itemize}
            \item 경쟁력 있는 이자율을 지불해야 함
        \end{itemize} &
        \begin{itemize}
            \item 상황이 나쁠 때 빠르게 사라지는 경향이 있음 (Runs away)
        \end{itemize} \\
        \addlinespace
        \textbf{은행 자본 (Bank Capital)} &
        \begin{itemize}
            \item 주식은 가격 위험에 노출되어 있어 예금보다 조달 비용이 비쌈
        \end{itemize} &
        \begin{itemize}
            \item 만기가 없어 가장 안정적 (Never matures, as sticky as it gets)
        \end{itemize} \\
        \bottomrule
    \end{tabular}
    \end{adjustbox}
\end{table}

\newpage
\subsection{Lesson 2-2.2: 은행 대출 (Bank Lending)}
\addcontentsline{toc}{subsection}{Lesson 2-2.2: 은행 대출 (Bank Lending)}

이제 은행 대출로 넘어가겠습니다. 논의를 세 부분으로 나누겠습니다.
\begin{enumerate}
    \item \textbf{대출의 정의:} 대출이 무엇인지 설명합니다.
    \item \textbf{고객 유형:} 은행이 대출 측면에서 어떤 고객을 다루는지 살펴봅니다.
    \item \textbf{대출 과정의 메커니즘:} 대출이 어떻게 발생하는지 설명합니다.
\end{enumerate}

\subsubsection*{대출의 정의와 조건}
공식적으로 대출이란 무엇일까요? 대출은 \textbf{금융 상품(Financial Instrument)}이며, \textbf{부채(Debt)}입니다. 제가 은행에서 돈을 빌릴 때, 저는 미리 정해진 계약 조건에 따라 상환하겠다는 법적 구속력이 있는 약속을 하는 것입니다.

\begin{itemize}
    \item \textbf{이자율 (Interest Rate):} 변동 금리(Variable Rate)일 수도 있고 고정 금리(Fixed Rate)일 수도 있습니다.
    \item \textbf{담보 (Collateral):} 대출은 어떤 담보로 보증될 수 있습니다.
    \item \textbf{약정 (Covenants):} 계약서에 성과 관련 조항이 명시될 수 있습니다.
    \item \textbf{선순위 채무 (Senior Debt):} 파산 시 사업 자산에 대한 최우선권을 가집니다.
\end{itemize}

은행은 일반적으로 대출 사업을 두 가지 범주로 나눕니다.
\begin{enumerate}
    \item \textbf{소비자 신용 (Consumer Credits):}
    \begin{itemize}
        \item 주거용 부동산 담보 대출: 주택 담보 대출(Mortgages) 및 주택 담보 대출(Home Equity Loans)을 포함합니다.
        \item 신용 카드(Credit Cards), 자동차 대출(Auto Loans), 학자금 대출(Student Loans)도 포함됩니다.
    \end{itemize}
    \item \textbf{상업 및 산업 대출 (Commercial and Industrial Lending):}
    \begin{itemize}
        \item 기업에 대한 대출을 포함합니다.
        \item 상업용 부동산 대출(Commercial Real Estate Loans), 투자 목적의 장기 대출인 기간 대출(Term Loans), 기업의 운전자금(Working Capital)으로 사용될 수 있는 사업 신용 한도(Business Lines of Credit) 등이 있습니다.
    \end{itemize}
\end{enumerate}

\subsubsection*{상업 은행의 고객 기반}
은행은 누구에게 대출을 해줄까요? 상업 은행의 고객 기반을 살펴보겠습니다.

\begin{enumerate}
    \item \textbf{대기업 (Larger Businesses):}
    \begin{itemize}
        \item 상업 은행을 우회하여 금융 시장에서 직접 자금을 조달하는 경향이 있습니다. 채권을 발행하고 주식 자금을 조달합니다.
        \item 이러한 기업에 대한 공개 정보가 많고 정보 비용이 적습니다.
        \item 때로는 대기업이 특별한 상황(예: 인수 거래 또는 부실 채무 구조 조정)에 있을 때 은행을 통해 대출을 주선하기도 합니다.
    \end{itemize}

    \item \textbf{중견 기업 대출 (Middle Market Lending):}
    \begin{itemize}
        \item 상업 은행 사업의 핵심입니다. 때로는 \textbf{관계형 대출(Relationship Lending)}이라고도 불립니다.
        \item 공개 정보가 제한적인 중견 기업들입니다. 아마도 공개 증권 거래소에 상장되어 있지 않을 것이며, 투자자들은 역선택(Adverse Selection)을 우려할 수 있습니다.
        \item 사업 고객을 찾고, 신용 분석을 수행하고, 새로운 대출에 자금을 조달하고, 관계를 유지하는 것이 대부분 은행의 대출 담당자(Loan Officers)와 신용 위원회(Credit Committees)의 핵심 역할입니다.
    \end{itemize}

    \item \textbf{소규모 및 신생 기업 (Small and New Businesses):}
    \begin{itemize}
        \item 일반적으로 어떤 형태의 보증(Guarantee) 없이는 은행 신용에 접근하기 어렵습니다.
        \item 사업주 개인의 보증, 중소기업청(Small Business Administration, SBA)의 부분 보증, 또는 사업주가 개인 자산(주택, 보트, 기타 금융 자산)을 담보로 제공할 수 있습니다.
        \item 소규모 기업은 위험하고 종종 실패합니다. 보증과 담보는 위험을 줄이고 은행이 대출에 참여하도록 돕습니다.
    \end{itemize}

    \item \textbf{개인 대출 (Loans to Individuals):}
    \begin{itemize}
        \item \textbf{주택 담보 대출 (Mortgage Credit):} 커뮤니티 은행과 지역 은행의 핵심 초점입니다. 이러한 대출은 정보 민감성이 높으며 차입자의 개별 상황에 따라 달라질 수 있습니다.
        \item \textbf{신용 카드, 학자금 대출, 자동차 대출:} 건당 대출 규모가 작고 비용이 많이 듭니다. 이 시장에는 큰 규모의 경제(Economies of Scale)가 존재하므로, 이러한 대출을 대규모로 다각화된 풀(Pool)로 만들 수 있는 대형 은행과 비은행 기관이 이 분야를 지배하는 경향이 있습니다.
    \end{itemize}
\end{enumerate}

\begin{figure}[h!]
    \centering
    \includegraphics[width=0.8\textwidth]{example_jpmorgan_segments.png} % 가상의 이미지 파일명
    \caption{JP모건 체이스의 사업 부문별 고객 분류}
    \label{fig:jpmorgan_segments}
    \textit{JP모건 체이스의 사업 부문별 고객 분류를 통해 이러한 다양한 고객들이 어떻게 나타나는지 볼 수 있습니다.}
\end{figure}

\begin{itemize}
    \item \textbf{소비자 및 커뮤니티 금융 (Consumer and Community Banking):} 제가 \textbf{소매 금융(Retail Banking)}이라고 부르는 부분입니다.
    \begin{itemize}
        \item \textbf{소비자 및 사업 금융 (Consumer and Business Banking):} 지점 네트워크를 통해 예금 수취, 소비자 신용 및 소규모 사업 대출을 제공합니다.
        \item \textbf{주택 대출 (Home Lending):} 주거용 주택 담보 대출 및 주택 담보 대출을 담당합니다.
        \item \textbf{카드 및 자동차 (Card and Auto):} 신용 카드와 자동차 대출 및 리스를 처리합니다.
        \item 이 모든 것은 개인 및 소규모 기업을 위한 것입니다.
    \end{itemize}

    \item \textbf{상업 금융 (Commercial Banking):} \textbf{도매 금융(Wholesale Banking)} 사업 내에 있습니다.
    \begin{itemize}
        \item \textbf{상업용 부동산 (Commercial Real Estate):} 사무 공간, 부동산 개발, 산업 공간 등에 대한 자금 조달을 포함합니다.
        \item \textbf{사업 금융 (Business Banking):} JP모건은 고객을 두 그룹으로 나눕니다.
        \begin{itemize}
            \item \textbf{중견 기업 (Middle-Market):} 연간 매출액이 2억 달러에서 5억 달러 사이의 중소기업을 대상으로 합니다. 이보다 작으면 커뮤니티 뱅킹 범주에 속합니다.
            \item \textbf{기업 고객 (Corporate Clients):} 연간 매출액이 5억 달러에서 20억 달러 사이의 대기업을 대상으로 합니다. 이보다 크면 투자 은행 부문에서 처리합니다.
        \end{itemize}
    \end{itemize}
\end{itemize}

\begin{figure}[h!]
    \centering
    \includegraphics[width=0.8\textwidth]{example_jpmorgan_diagram.png} % 가상의 이미지 파일명
    \caption{JP모건의 비즈니스 모델: 도매 및 소매 금융 매핑}
    \label{fig:jpmorgan_diagram}
    \textit{JP모건의 비즈니스 모델을 이전 다이어그램에 매핑할 수 있습니다. 여기는 중견 기업에 대한 대출을 다루는 도매 사업이고, 여기는 소규모 기업과 개인에 대한 대출을 다루는 소매 사업입니다.}
\end{figure}

\begin{summarybox}{Lesson 2-2.2 요약}
이것이 상업 은행 비즈니스 모델의 개요입니다. 은행은 다양한 출처로부터 자금을 조달하고 광범위한 대출을 제공합니다.
\end{summarybox}

\newpage
\subsection{Lesson 2-2.3: 대출 과정의 변화 (Changes in the Lending Process)}
\addcontentsline{toc}{subsection}{Lesson 2-2.3: 대출 과정의 변화 (Changes in the Lending Process)}

논의를 마무리하기 전에 두 가지를 더 다루고자 합니다. 첫째, 대출 과정 자체에 대해 좀 더 정확하게 설명하고 싶습니다. 둘째, 지난 50년 동안 이 대출 과정이 어떻게 변했는지 설명하고 싶습니다.

\subsubsection*{대출 과정의 세 가지 필수 단계}
대출 과정에는 세 가지 필수 단계가 있습니다.
\begin{enumerate}
    \item \textbf{개시 (Origination):}
    \begin{itemize}
        \item 대출을 승인하는 단계입니다. 대출 담당자(Loan Officer)가 고객과 만나 신용도를 판단합니다.
        \item 이는 차입자의 재정 상황이나 사업 계획에 대한 신중한 분석을 포함합니다.
        \item 대출 담당자는 대출 금액, 만기, 이자율 및 수수료, 담보 등 조건을 협상합니다.
        \item 우리는 이 단계를 이전 논의에서 \textbf{스크리닝(Screening)}이라고 불렀습니다.
    \end{itemize}

    \item \textbf{자금 조달 (Financing):}
    \begin{itemize}
        \item 대출을 약정한 후, 은행은 차입자에게 실제로 제공할 자금을 마련해야 합니다.
    \end{itemize}

    \item \textbf{관리 (Servicing):}
    \begin{itemize}
        \item 이 단계는 대출이 이루어진 후 은행이 차입자와의 관계를 지속하는 것을 포함합니다.
        \item 은행은 차입자로부터 상환금을 수취합니다.
        \item 은행은 차입자의 진행 상황과 재정 상태를 모니터링하여 은행이 전액 상환받을 가능성을 극대화합니다.
    \end{itemize}
\end{enumerate}

\subsubsection*{전통적인 모델: 개시 및 보유 (Originate and Hold)}
전통적으로 은행은 이 세 단계를 모두 수행했습니다.
\begin{itemize}
    \item 대출을 승인하고,
    \item 예금을 사용하여 대출에 자금을 조달하며,
    \item 대출이 만기될 때까지 장부에 보유하면서 상환금을 수취했습니다.
\end{itemize}
이것이 우리가 \textbf{개시 및 보유 모델(Originate and Hold Model)}이라고 부르는 것입니다.

\subsubsection*{현대적인 모델: 개시 및 분배 (Originate and Distribute)}
이제 재미있는 부분입니다. 현대 시대에는 이 세 단계 각각이 \textbf{개시 및 분배 모델(Originate and Distribute Model)}이라고 불리는 방식으로 분리될 수 있습니다. 이 비즈니스 모델 하에서 상업 은행은 \textbf{수수료 수입(Fee Income)}에 중점을 둡니다.

\begin{itemize}
    \item \textbf{개시 (Origination):}
    \begin{itemize}
        \item 은행은 여전히 차입자와 만나 대출을 개시합니다. 이 부분은 변하지 않았지만, 때로는 브로커(Broker)의 도움을 받기도 합니다.
        \item 주택 담보 대출 시장에서는 온라인 브로커가 수수료를 받고 대출 기관과 고객을 연결해주는 경우를 접했을 수 있습니다.
    \end{itemize}

    \item \textbf{자금 조달 (Financing):}
    \begin{itemize}
        \item 다른 금융기관이 주택 구매 자금을 제공합니다. 이는 보험 회사와 같은 투자자일 수도 있고, \textbf{증권화(Securitization)}라는 추가 단계가 있을 수도 있습니다.
    \end{itemize}

    \item \textbf{관리 (Servicing):}
    \begin{itemize}
        \item 은행이 관리를 수행할 수도 있지만, 이마저도 때로는 아웃소싱됩니다.
        \item 매월 상환금을 은행에 지불하는 대신, 주택 담보 대출 관리자(Mortgage Servicer)에게 지불하고, 관리자는 그 상환금을 투자자에게 전달합니다.
    \end{itemize}
\end{itemize}

\begin{center}
    \textit{왜 이렇게 할까요?}
\end{center}

일부 은행이 정보 비용을 최소화하고 최적의 차입자를 승인하는 데 정말 뛰어나다면, \textbf{비교 우위(Comparative Advantage)} 이론은 그들이 그 활동에 집중해야 한다고 말합니다. 그리고 다른 파트너가 자금을 제공할 수 있습니다. 이것이 바로 \textbf{개시 및 분배 모델(Originate and Distribute Model)}의 핵심입니다.

\subsubsection*{핀테크(Fintech)의 등장과 대출 시장의 변화}
마지막으로 한 가지 더 생각할 점이 있습니다. 이 이론은 21세기에 접어들면서 무너질 수도 있습니다. 요즘에는 로켓 모기지(Rocket Mortgage)와 같은 \textbf{핀테크(Fintech)} 기업들이 주택 담보 대출 시장을 뒤흔들었습니다. 소비자로서 제가 주택 담보 대출을 재융자하거나 특정 품질 기준을 충족하는 우량 차입자라면, 온라인 대출 기관으로부터 빠르고 편리하게 대출을 받을 수 있습니다. 이런 상황에서 은행이 비교 우위를 가질 수 있을까요?

우리는 아마존(Amazon)과 같은 다른 기술 기업들에 의한 기업 대출 시장에서도 유사한 진입을 보고 있습니다. 이들 기업은 고객의 관심과 사업 매출에 대한 방대한 양의 정보를 실시간으로 수집합니다. 아마존이 어떤 차입자에게 대출을 승인해야 할지 결정하는 데 더 나은 위치에 있을 수도 있습니다. 우리는 지켜봐야 할 것입니다.

\begin{summarybox}{Lesson 2-2.3 요약}
이번 세션에서는 상업 은행 사업을 면밀히 살펴보았습니다.
\begin{itemize}
    \item 은행이 어떻게 자금을 조달하고 대출하는지 설명했습니다.
    \item 소매 금융과 도매 금융의 차이점을 설명했습니다.
    \item 마지막으로, 은행이 최적의 차입자를 선택하는 데 비교 우위를 가진다면 합리적인 \textbf{개시 및 분배 모델(Originate and Distribute Model)}을 소개했습니다.
\end{itemize}
\end{summarybox}

\newpage

\newpage
\section*{개요}
이 문서는 'Banking and Financial Institutions Module 2'의 핵심 내용을 다루며, 투자은행의 조직 구조, 주요 활동, 그리고 현대 은행의 재무제표 분석에 초점을 맞춥니다. 투자은행이 기업과 금융시장을 연결하는 방식, 주요 수익원, 그리고 상업은행과의 통합 과정을 이해합니다. 또한, 상업은행의 대차대조표(자산 및 부채)를 상세히 살펴보고, 대차대조표 외 활동(Off-Balance Sheet Activities)의 중요성과 위험성을 분석하여 금융기관의 복잡한 운영 방식을 종합적으로 파악하는 것을 목표로 합니다.

\newpage
\section*{용어 정리}

\begin{adjustbox}{width=\textwidth,center}
\begin{tabular}{llll}
\toprule
\textbf{용어} & \textbf{쉬운 설명} & \textbf{원어} & \textbf{비고} \\
\midrule
투자은행 & 기업이 주식이나 채권을 발행하여 자금을 조달하도록 돕는 금융기관 & Investment Bank & 상업은행과 달리 대출보다는 자본시장 접근 지원 \\
상업은행 & 예금을 받아 대출을 해주는 전통적인 은행 & Commercial Bank & 일반 개인 및 기업 대상 \\
고액 자산가 & 상당한 금융 자산을 보유한 개인 & High-Net-Worth Individual & 투자은행의 주요 고객 중 하나 \\
증권 인수 & 기업이 발행하는 주식이나 채권을 투자은행이 매입하여 시장에 판매하는 과정 & Underwriting & 투자은행의 핵심 수익원 \\
증권 거래 & 주식, 채권 등 증권을 사고파는 활동 & Securities Trading & 고객 주문 처리 및 자기 계정 거래 포함 \\
M\&A 자문 & 기업의 인수합병(Mergers and Acquisitions) 또는 구조조정에 대한 자문 서비스 & M\&A Advisory & 기업 전략 및 가치 평가 전문성 요구 \\
신디케이트론 & 여러 은행이 공동으로 대규모 대출을 제공하여 위험을 분산하는 방식 & Syndicated Loan & 주로 대규모 기업 대출에 활용 \\
북빌딩 & 투자은행이 신규 증권 발행 시 투자자들의 수요를 파악하고 가격을 결정하는 과정 & Book-building & IPO 과정의 핵심 단계 \\
비드-애스크 스프레드 & 매수호가(Bid Price)와 매도호가(Ask Price)의 차이 & Bid-Ask Spread & 시장 조성자(Market Maker)의 주요 수익원 \\
시장 조성 & 투자자들이 증권을 쉽게 사고팔 수 있도록 매수호가와 매도호가를 제시하는 활동 & Market-making & 유동성 공급 역할 \\
자기 계정 거래 & 투자은행이 고객을 위해서가 아니라 자체 이익을 위해 증권을 사고파는 활동 & Proprietary Trading & 높은 수익 가능성, 높은 위험 수반 \\
대차대조표 & 특정 시점의 기업 자산, 부채, 자본 상태를 보여주는 재무 보고서 & Balance Sheet & 은행의 재무 건전성 파악에 필수적 \\
뱅킹 북 & 전통적인 대출 및 예금 활동과 관련된 은행의 자산 및 부채 & Banking Book & 장기적인 이자 수익 및 비용 발생 \\
트레이딩 북 & 단기적인 증권, 상품, 통화 거래와 관련된 은행의 자산 및 부채 & Trading Book & 단기 매매 차익을 목표로 함 \\
대차대조표 외 활동 & 현재는 자산이나 부채로 기록되지 않지만, 특정 사건 발생 시 재무제표에 영향을 미칠 수 있는 활동 & Off-Balance Sheet (OBS) Activities & 신용카드 약정, 파생상품 등이 대표적 \\
우발 자산/부채 & 특정 미래 사건의 발생 여부에 따라 자산이나 부채가 될 수 있는 항목 & Contingent Asset/Liability & OBS 활동의 핵심 개념 \\
미사용 대출 약정 & 은행이 고객에게 약속했지만 아직 인출되지 않은 대출 한도 & Undrawn Loan Commitment & 신용카드 한도 등이 해당 \\
파생상품 & 기초자산의 가치 변동에 따라 가치가 결정되는 금융상품 & Derivatives & 위험 헤지 또는 투기 목적으로 사용 \\
명목 금액 & 파생상품 계약의 기초자산 총액 & Notional Amount & 실제 시장 가치와는 다름 \\
신용부도스왑 & 채권이나 대출의 부도 위험에 대한 보험 계약 & Credit Default Swap (CDS) & 신용 위험을 전가하는 수단 \\
\bottomrule
\end{tabular}
\end{adjustbox}

\newpage
\section{투자은행의 조직 구조와 활동}

\subsection{투자은행이란 무엇인가?}
투자은행은 기업이나 다른 기관들이 금융시장에 접근할 수 있도록 돕는 금융기관입니다.
\begin{itemize}
    \item \textbf{상업은행과의 차이:} 일반적인 기업이 대출을 원하면 상업은행(Commercial Bank)에 가지만, 주식이나 채권을 발행하여 자금을 조달하고 싶다면 투자은행(Investment Bank)을 찾습니다.
    \item \textbf{고액 자산가 서비스:} 고액 자산가(High-Net-Worth Individual)가 주식 거래를 실행하고 싶을 때도 투자은행가에게 연락할 수 있습니다.
\end{itemize}
이러한 구분은 단순화된 설명이지만, 투자은행의 핵심 역할을 이해하는 데 도움이 됩니다. 투자은행의 목표는 차입자(Borrower)와 저축자(Saver)를 금융시장에서 직접 연결하는 것입니다.

\begin{summarybox}
\textbf{핵심 요약:} 투자은행은 기업의 자금 조달(주식/채권 발행) 및 M\&A를 돕고, 증권 거래 서비스를 제공하는 금융기관입니다. 상업은행이 대출과 예금에 집중하는 것과 달리, 투자은행은 자본시장을 통한 자금 조달과 투자에 특화되어 있습니다.
\end{summarybox}

\subsection{투자은행의 세 가지 주요 사업 분야}
투자은행은 크게 세 가지 주요 사업 라인을 가지고 있습니다. 이 활동들은 모두 수수료를 발생시키며, 상업은행과 달리 장기 자금 조달이나 광범위한 지점 네트워크가 필요하지 않아 매우 수익성이 높은 사업이 될 수 있습니다.

\begin{enumerate}
    \item \textbf{자금 조달 (Financing):}
    \begin{itemize}
        \item \textbf{정의:} 투자은행은 고객(주로 대기업 및 정부)을 대신하여 증권을 발행(Originate), 인수(Underwrite), 그리고 금융시장에 배치(Place)합니다.
        \item \textbf{예시:} 기업이 새로운 공장을 짓기 위해 자금이 필요할 때, 투자은행은 해당 기업의 주식이나 채권을 발행하여 투자자들에게 판매하는 과정을 총괄합니다.
        \item \textbf{수익원:} 증권 인수 수수료(Underwriting Fees)가 발생합니다.
    \end{itemize}

    \item \textbf{증권 거래 (Securities Trading):}
    \begin{itemize}
        \item \textbf{정의:} 증권 발행 및 판매 사업에 이미 참여하고 있기 때문에, 대부분의 투자은행은 증권 거래에도 관여합니다. 고객의 매수 및 매도 주문을 실행하고, 종종 자체 이익을 위해 거래하기도 합니다.
        \item \textbf{예시:} 고객이 특정 주식을 사고 싶을 때, 투자은행은 고객을 대신하여 시장에서 주식을 매수하고 수수료를 받습니다. 또한, 은행 자체적으로 주식이나 채권을 사고팔아 이익을 얻기도 합니다 (자기 계정 거래).
        \item \textbf{수익원:} 거래 실행에 대한 수수료(Commissions)가 발생합니다.
    \end{itemize}

    \item \textbf{M\&A 및 기업 구조조정 자문 (M\&A Advisory and Corporate Restructuring):}
    \begin{itemize}
        \item \textbf{정의:} 투자은행은 기업의 인수합병(Mergers and Acquisitions, M\&A) 또는 기타 기업 구조조정 활동을 촉진하는 데 도움을 줍니다.
        \item \textbf{예시:} 한 기업이 다른 기업을 인수하려고 할 때, 투자은행은 인수 대상 기업의 가치를 평가하고, 협상을 돕고, 필요한 자금 조달 방안을 제시합니다.
        \item \textbf{수익원:} 자문 서비스에 대한 수수료(Advisory Fees)가 발생합니다.
    \end{itemize}
\end{enumerate}

\subsection{투자은행의 기업 구조: JP Morgan Chase 사례}
JP Morgan Chase의 사례를 통해 투자은행의 기업 구조를 살펴보겠습니다.
\begin{itemize}
    \item \textbf{도매 사업 부문 (Wholesale Business Arm):} JP Morgan의 기업 및 투자은행 부문은 은행의 도매 사업 부문 내에 위치합니다. 이는 일반 소매 고객(Retail Customers)이 아닌 대규모 고객(Big Clients)을 주로 상대한다는 의미입니다.
    \item \textbf{두 가지 핵심 사업:} 투자은행은 크게 '뱅킹(Banking)'과 '마켓(Markets)' 두 가지 사업으로 구성됩니다.
    \begin{itemize}
        \item \textbf{뱅킹 (Banking):} 기업 전략 및 구조 자문, 자본 조달(주식, 채권 시장), 신디케이트론 시장 자문 등을 포함합니다.
        \item \textbf{마켓 (Markets):} 현금 및 파생상품 거래, 중개 서비스, M\&A 자문, 증권 인수, 증권 거래 등을 포괄합니다.
    \end{itemize}
\end{itemize}
JP Morgan은 2020년에 약 85억 달러의 투자은행 수수료를 벌어들였는데, 이는 연간 수익의 상당 부분을 차지합니다.

\subsection{투자은행 산업의 분류}
투자은행 산업에는 다양한 전문 기업들이 존재하지만, 가장 큰 풀서비스 투자은행은 세 가지 유형으로 분류할 수 있습니다.

\begin{enumerate}
    \item \textbf{상업은행 지주회사의 자회사 (예: JP Morgan):}
    \begin{itemize}
        \item \textbf{특징:} JP Morgan의 투자은행 부문은 상업은행 지주회사(Commercial Bank Holding Company)의 자회사입니다.
        \item \textbf{고객 기반:} 기업, 정부, 금융기관, 고액 자산가 등 광범위한 고객층을 보유한 가장 큰 투자은행들입니다.
    \end{itemize}

    \item \textbf{특정 고객 기반에 특화된 투자은행 (예: Goldman Sachs):}
    \begin{itemize}
        \item \textbf{특징:} Goldman Sachs와 같은 투자은행은 주로 기업 및 증권 거래 활동이 활발한 고객에 집중하는 경향이 있습니다.
        \item \textbf{고객 기반:} 보다 전문화된 고객 기반을 가집니다.
    \end{itemize}

    \item \textbf{독립형 투자은행 (Standalone Investment Banks, 예: Lazard):}
    \begin{itemize}
        \item \textbf{특징:} Lazard와 같은 독립형 투자은행은 모든 범위의 투자은행 서비스를 제공하지만, 지리적 범위가 제한적입니다. 주로 주요 도시에만 운영되며, 대규모 고객에 집중합니다.
        \item \textbf{역사적 변화:} 과거에는 Merrill Lynch, Bear Stearns, Lehman Brothers, Morgan Stanley 등 여러 독립형 투자은행이 있었으나, 규제 변화와 2008년 금융 위기로 인해 대부분 상업은행과 합병되거나 상업은행으로 재편되었습니다.
    \end{itemize}
\end{enumerate}

\begin{tipbox}
\textbf{독립형 vs. 계열사:} '독립형(Standalone)'이란 상업은행과 제휴하지 않은 투자은행을 의미합니다. 과거에는 독립형이 많았으나, 1990년대 후반의 규제 완화(상업은행과 투자은행 활동 겸업 허용)와 2008년 금융 위기(대규모 인수합병 및 재편)로 인해 현재는 대부분 상업은행 지주회사 산하에 편입되어 있습니다.
\end{tipbox}

\subsection{투자은행 산업의 변화와 통합}
투자은행 산업은 큰 변화를 겪었습니다.
\begin{itemize}
    \item \textbf{규제 변화 (1990년대 후반):} 1990년대 후반의 규제 변화로 상업은행과 투자은행 활동이 같은 회사에서 가능해지면서, 여러 상업은행이 자체 투자은행을 인수하거나 설립했습니다. Gramm-Leach-Bliley Act 통과와 함께 Citigroup이 Citicorp(상업은행)와 Travelers(투자은행을 소유한 보험회사)의 합병으로 탄생한 것이 대표적인 예입니다.
    \item \textbf{금융 위기 (2008년):} 2008년 금융 위기는 대규모 부실 인수(Distressed Acquisitions)의 물결을 일으켰습니다.
    \begin{itemize}
        \item Lehman Brothers는 파산했습니다.
        \item Bear Stearns와 Merrill Lynch는 상업은행에 인수되었습니다 (Bank of America가 Merrill Lynch를 인수한 것이 가장 큰 부실 인수였습니다).
        \item Goldman Sachs와 Morgan Stanley는 생존을 위해 상업은행으로 재편해야 했습니다.
    \end{itemize}
\end{itemize}

\begin{table}[h!]
\centering
\caption{미국 투자은행 산업의 주요 합병 사례}
\label{tab:us_ib_mergers}
\begin{adjustbox}{width=\textwidth,center}
\begin{tabular}{lll}
\toprule
\textbf{합병 주체} & \textbf{피합병 주체} & \textbf{합병 연도} \\
\midrule
Citigroup & Citicorp (상업은행) + Travelers (보험/투자은행) & 1998 \\
Chase & JPMorgan & 직후 \\
Bank of America & Merrill Lynch & 2008 \\
\bottomrule
\end{tabular}
\end{adjustbox}
\end{table}

\subsection{글로벌 투자은행 수익 현황 (2020년)}
2020년 투자은행 수익 분석에 따르면, 전 세계 상위 10개 투자은행 중 Jefferies만이 독립형 회사이며, 나머지는 모두 상업은행과 제휴하여 소매 고객으로부터 예금을 받습니다.
\begin{itemize}
    \item \textbf{시장 집중도:} 이 상위 10개 은행이 전 세계 수수료의 약 50\%를 차지하며, 시장 선두 주자는 약 10\%를 차지합니다.
    \item \textbf{경쟁 환경:} 상당한 통합이 있었음에도 불구하고, 글로벌 규모에서는 여전히 상당히 경쟁적인 시장입니다.
\end{itemize}

\newpage
\section{투자은행 활동의 세부 사항}

\subsection{증권 인수 (Securities Underwriting)}
증권 인수는 투자은행 수수료의 대부분을 차지하는 핵심 활동입니다. 주로 주식, 채권, 신디케이트론 인수에 초점을 맞춥니다.

\begin{enumerate}
    \item \textbf{과정:}
    \begin{itemize}
        \item \textbf{의뢰 (Mandate):} 기업이 5억 달러와 같은 대규모 자금을 조달하기 위해 금융 상품(주식 또는 채권)을 판매하고자 할 때, 여러 투자은행이 이 인수 과정을 주도하기 위한 의뢰(Mandate)를 따내기 위해 경쟁합니다.
        \item \textbf{주간사 (Lead Underwriter):} 주간사로 선정된 투자은행은 거래 조건을 설정하고 새로 발행된 금융 상품의 판매를 관리하는 역할을 합니다.
    \end{itemize}

    \item \textbf{주식 발행 (Equity Sell):}
    \begin{itemize}
        \item \textbf{IPO (Initial Public Offering):} 기업이 처음으로 주식을 판매하는 경우를 기업공개(IPO)라고 합니다.
        \item \textbf{추가 발행 (Follow-on Deal/Seasoned Offering):} 이미 상장된 기업이 추가로 주식을 판매하는 경우입니다.
        \item \textbf{북빌딩 (Book-building):} 투자은행은 기업 및 잠재 투자자들과 만나 주식 가격과 매수 의사를 협상합니다. 이 과정을 북빌딩이라고 합니다.
        \item \textbf{인수 및 판매:} 투자은행은 전체 발행 물량을 보장된 가격에 매입하여 외부 투자자에게 더 높은 가격에 판매할 수 있습니다. 매입 가격과 판매 가격의 차이(Spread)가 투자은행의 이익이 됩니다.
        \item \textbf{예시:} 투자은행이 1억 주를 주당 5달러에 매입하여 주당 5.25달러에 판매하면, 주당 0.25달러의 스프레드로 2,500만 달러의 이익을 얻습니다.
        \item \textbf{위험:} 그러나 이러한 이익은 보장되지 않습니다. Morgan Stanley는 Facebook IPO 가격을 잘못 책정하여 주가가 몇 주 만에 거의 50\% 하락하는 사태를 겪었습니다. 이는 투자은행의 명성과 향후 사업 전망에 부정적인 영향을 미칠 수 있습니다. 따라서 인수 가격은 재고를 이익을 내고 소진할 수 있으면서도, 기업의 공정한 가치를 반영해야 합니다.
    \end{itemize}

    \item \textbf{채권 인수 (Bond Underwriting):} 주식 인수와 유사한 과정을 따릅니다.

    \item \textbf{신디케이트론 (Syndicated Loan):}
    \begin{itemize}
        \item \textbf{정의:} 일반적으로 상업은행이 기업 대출을 제공하지만, 대규모 대출의 경우 단일 은행이 감당하기에는 위험이 너무 큽니다. 이때 여러 은행과 대출 기관이 모여 공동으로 대출을 제공하여 위험을 분산합니다.
        \item \textbf{주간사 역할:} 투자은행(인수 은행)은 대출 신디케이트를 대표하여 대출 조건을 협상합니다. 대출 과정의 시작(Origination)과 사후 관리(Servicing)를 모두 담당합니다.
        \item \textbf{수익원:} 신디케이트를 주도하는 은행은 많은 수수료를 얻습니다. 다른 신디케이트 구성원들은 자금을 제공하는 '침묵하는 파트너' 역할을 하며, 대출 만기까지 일반적인 이자 수입을 얻습니다.
    \end{itemize}
\end{enumerate}

\subsection{M\&A 자문 (M\&A Advisory)}
투자은행은 기업의 인수합병(M\&A) 거래를 협상하고 실행하는 데 자주 관여합니다.

\begin{itemize}
    \item \textbf{역할:}
    \begin{itemize}
        \item \textbf{인수 대상 탐색:} 기업 X가 기업 Y를 인수하고 싶지만 적절한 가격을 모를 때, 투자은행은 잠재적 인수 대상에 대한 조사를 수행합니다.
        \item \textbf{매수자 탐색:} 기업 Z가 사업 부문을 매각하고 싶을 때, 투자은행은 잠재적 매수자를 찾는 데 도움을 줍니다.
        \item \textbf{가치 평가:} 재무 모델링 기법을 사용하여 상장 및 비상장 기업의 가치를 추정하는 전문가입니다.
        \item \textbf{자금 조달 지원:} M\&A 자문 은행은 거래 자금 조달을 위한 새로운 증권 인수(주식, 채권, 신디케이트론 등 모든 옵션 가능)도 지원할 수 있습니다.
    \end{itemize}
    \item \textbf{보상:} M\&A 거래의 세부 사항을 확정하는 것은 재정적, 법률적 관점에서 복잡하기 때문에, 투자은행은 제공하는 자문 및 서비스에 대해 높은 보상을 받습니다.
\end{itemize}

\subsection{투자은행 상품별 수수료 수익 (2020년)}
2020년 투자은행의 다양한 상품별 수수료 수익을 살펴보면 다음과 같습니다.
\begin{itemize}
    \item \textbf{주식 인수:} IPO(Initial Public Offering)와 추가 발행(Follow-on Offerings)으로 나뉩니다.
    \item \textbf{채권 인수:} 투자 등급(Investment Grade, 신용 등급이 좋은 기업)과 투자 등급 미만(Below Investment Grade, 신용 등급이 낮은 기업)으로 나뉩니다.
    \item \textbf{수익 규모:} 2020년에는 주식 인수, 채권 인수, M\&A 관련 활동에서 발생하는 총 수수료 수익이 모두 비슷한 규모였습니다. 신디케이트론에서 발생하는 수수료는 이들의 약 절반 수준이었습니다.
\end{itemize}

\newpage
\subsection{증권 거래 (Securities Trading)}
증권 거래는 투자은행의 또 다른 핵심 활동이며, 시장 조성(Market-making)과 자기 계정 거래(Proprietary Trading)로 구성됩니다.

\subsubsection{시장 조성 (Market-making)}
\begin{itemize}
    \item \textbf{정의:} 투자은행이 투자자들이 증권을 거래할 수 있는 시장을 만드는 활동입니다. 많은 증권 회사들이 이 역할을 하지만, 투자은행이 중심적인 역할을 합니다.
    \item \textbf{브로커 역할:} 브로커는 고객을 대신하여 거래를 실행하고 수수료나 커미션을 받습니다.
    \item \textbf{딜러 역할:} 투자은행의 딜러 부문은 자체적으로 재고(Inventory)로 일부 주식을 보유할 수 있습니다. 이 경우 고객은 은행 자체로부터 주식을 매수할 수도 있습니다.
    \item \textbf{비드-애스크 스프레드 (Bid-Ask Spread):} 투자은행이 Facebook 주식을 매도호가(Ask Price) 20달러에 매수하여 고객에게 매수호가(Bid Price) 21달러에 판매한다면, 1달러의 이익을 얻습니다. 이 이익을 비드-애스크 스프레드라고 합니다.
    \item \textbf{시너지 효과:} 증권 딜링과 인수 활동을 결합하면 명확한 시너지 효과가 있습니다. 새로운 주식을 인수한 투자은행은 남은 주식을 재고로 보유하기 쉽습니다.
\end{itemize}

\subsubsection{자기 계정 거래 (Proprietary Trading)}
\begin{itemize}
    \item \textbf{정의:} 은행이 고객을 돕는 것이 아니라, 자체 이익을 위해 금융 상품에 적극적인 포지션을 취하는 활동입니다.
    \item \textbf{포지션 거래 (Position Trade):} 유리한 가격 변동을 예상하여 몇 주 또는 몇 달 동안 증권에 대해 롱(Long) 또는 숏(Short) 포지션을 취하는 것입니다.
    \item \textbf{차익 거래 (Arbitrage Trading):} 금융시장의 작고 일시적인 가격 불일치(Mispricing)를 이용하는 것입니다. 이론적으로는 존재해서는 안 되는 이러한 불일치는 이익을 얻을 기회를 제공하며, 시장의 유동성을 높이고 가격 효율성을 개선하는 데 기여하기도 합니다.
\end{itemize}

\subsection{투자은행 및 거래 수익 (2020년)}
미국 및 유럽의 12개 대형 금융기관의 총 투자은행 및 거래 수익은 2020년에 약 2,000억 달러에 육박했습니다.
\begin{itemize}
    \item \textbf{수익 구성:} 이 중 약 1,500억 달러(약 75\%)가 채권 및 주식 거래에서 발생했습니다.
    \item \textbf{수수료 수익:} 나머지 약 25\%는 투자은행 수수료에서 발생했습니다.
\end{itemize}
이는 거래 활동이 투자은행의 전체 수익에서 매우 큰 비중을 차지함을 보여줍니다.

\begin{summarybox}
\textbf{투자은행의 주요 활동 요약:}
\begin{itemize}
    \item \textbf{증권 인수 (Securities Underwriting):} 기업이 새로운 증권을 발행하고 투자자에게 판매하도록 돕습니다.
    \item \textbf{거래 (Trading):} 이미 발행된 증권의 거래를 돕고, 자체적으로 거래하여 이익을 창출합니다.
    \textbf{기업 전략 자문 (Corporate Strategy):} M\&A 및 기업 구조조정에 대한 자문을 제공합니다.
\end{itemize}
이 세 가지 활동은 모두 투자은행의 중요한 수수료 기반 수입원입니다.
\end{summarybox}

\newpage
\section{현대 은행의 재무제표 분석}

\subsection{대차대조표 (Balance Sheet)}
은행이 기업으로서 어떻게 조직되고, 은행 산업 및 비즈니스 모델이 어떻게 진화했는지 이해했습니다. 이제 은행이 무엇을 하는지 더 명확히 이해하기 위해 상업은행 지주회사의 대차대조표를 자세히 살펴보겠습니다. 대차대조표는 자산(Assets) 측면에서 투자 목록을, 부채(Liabilities) 측면에서 자금 조달원을 보여줍니다. JP Morgan Chase의 대차대조표와 전체 상업은행 산업의 통합 대차대조표를 사례 연구로 활용할 것입니다.

\begin{summarybox}
\textbf{핵심 요약:} 대차대조표는 특정 시점의 은행 재무 상태를 보여주는 스냅샷입니다. '자산'은 은행이 소유한 것(투자), '부채'는 은행이 빚진 것(자금 조달원), '자본'은 소유주의 몫을 나타냅니다.
\end{summarybox}

\subsubsection{상업은행 대차대조표의 기본 구조}
상업은행 대차대조표의 간단한 구조는 다음과 같습니다.

\begin{itemize}
    \item \textbf{자산 (Asset Side):} 은행이 조달한 자금을 어떻게 사용하는지를 보여줍니다.
    \begin{itemize}
        \item \textbf{현금 (Cash):} 금고에 있는 현금, 다른 은행의 예금 계좌.
        \item \textbf{대출 (Loans):} 주택 담보 대출, 기업 대출, 가계 대출 등.
        \item \textbf{기타 금융 자산 (Other Financial Assets):} 국채, 회사채, 기타 증권 등.
        \item \textbf{유형 자산 (Physical Assets):} 건물, 비품 등 (예: 대출 담당자가 앉는 테이블과 의자).
    \end{itemize}
    \item \textbf{부채 (Liability Side):} 은행이 어떻게 자금을 조달하는지를 보여줍니다.
    \begin{itemize}
        \item \textbf{고객 예금 (Customer Deposits):} 당좌 예금, 저축 예금 등.
        \item \textbf{도매 자금 조달 (Wholesale Funding):} 머니 마켓에서 조달하는 자금.
        \item \textbf{은행 자본 (Bank Capital/Equity):} 은행 소유주가 제공한 자금, 유보 이익(Retained Earnings), 추가 주식 발행 등.
    \end{itemize}
\end{itemize}

\subsubsection{뱅킹 북과 트레이딩 북}
대규모 상업은행 지주회사는 상당한 시장 조성 및 증권 거래 활동을 수행하므로, 대차대조표를 '뱅킹 북(Banking Book)'과 '트레이딩 북(Trading Book)'으로 구분합니다.

\begin{itemize}
    \item \textbf{뱅킹 북 (Banking Book):}
    \begin{itemize}
        \item \textbf{특징:} 전통적인 은행의 장기 차입 및 대출 활동을 나타냅니다.
        \item \textbf{수익원:} 대출에서 발생하는 이자 수입과 수수료를 창출합니다.
        \item \textbf{기록 방식:} 장기적인 관점에서 자산과 부채를 기록합니다.
    \end{itemize}
    \item \textbf{트레이딩 북 (Trading Book):}
    \begin{itemize}
        \item \textbf{특징:} 금융시장에서 빠르게 사고팔 수 있는 증권, 상품, 통화에 대한 투자를 포함합니다.
        \item \textbf{기록 방식:} 롱 포지션(매수 포지션)은 자산 측에 시장 가치로 기록되고, 숏 포지션(차입을 통한 매도 포지션)은 부채로 기록됩니다.
        \item \textbf{수익원:} 거래 포트폴리오에서 발생하는 거래 이익 및 손실은 비이자 수익 및 비용으로 기록됩니다.
    \end{itemize}
\end{itemize}

\subsubsection{JP Morgan Chase의 대차대조표 분석 (2020년)}
JP Morgan Chase는 세계에서 가장 크고 복잡한 은행 중 하나입니다. 2020년 소매 및 도매 사업 부문을 통합한 대차대조표를 살펴보겠습니다. 모든 숫자는 백만 달러 단위입니다.

\begin{itemize}
    \item \textbf{총 자산 (Total Assets):} 약 3.4조 달러 (Trillion).
    \item \textbf{총 부채 (Total Liabilities):} 약 3.1조 달러.
    \item \textbf{총 주주 자본 (Stockholder's Equity):} 약 3,000억 달러 (자산 - 부채).
\end{itemize}

\paragraph{자산 측면의 주요 항목:}
\begin{enumerate}
    \item \textbf{현금 (Cash):} 은행이 보유한 현금과 다른 은행의 예금 계좌에 예치된 현금을 포함합니다. 이는 은행이 보유한 가장 유동적인 자산입니다.
    \item \textbf{연방 기금 판매 및 환매 조건부 계약 (Federal Funds Sold and Reverse Repo Agreements):} JP Morgan이 다른 금융기관(예: 다른 은행)에 단기적으로 자금을 대출해 주는 것을 의미합니다.
    \item \textbf{거래 계정 자산 (Trading Account Assets):} 단기적인 포지션 또는 차익 거래의 일환으로 매도할 의도로 보유하는 증권입니다.
    \item \textbf{투자 증권 (Investment Securities):} 은행의 거래 포트폴리오에 속하지 않는 장기 채무 증권 투자입니다.
    \item \textbf{대출 (Loans):} 대차대조표에서 가장 큰 항목으로, 약 1조 달러에 달합니다.
    \begin{itemize}
        \item \textbf{소매 부문 대출:} 약 5,000억 달러는 개인 및 소기업 대출, 주거용 부동산 담보 대출, 신용카드, 자동차 대출 등 소매 부문에서 발생합니다.
        \item \textbf{도매 부문 대출:} 나머지 5,000억 달러는 대기업 및 정부에 대한 대출 등 도매 부문에서 발생합니다.
    \end{itemize}
\end{enumerate}
JP Morgan은 광범위한 지점 네트워크를 통해 약 1조 달러의 예금을 조달하며, 이는 소매 뱅킹 사업의 대출 규모보다 많습니다. 이는 JP Morgan과 다른 은행들이 상대적으로 저렴하고 안정적인 예금 자금을 사용하여 다른 투자 활동을 지원한다는 것을 의미합니다.

\paragraph{부채 측면의 주요 항목:}
\begin{enumerate}
    \item \textbf{예금 (Deposits):} JP Morgan의 가장 큰 자금 조달원입니다. 당좌 예금(Transaction Accounts)과 비거래 계좌(Non-Transaction Accounts, 주로 단기 채권과 유사한 예금 증서)를 포함합니다.
    \item \textbf{차입금 (Borrowed Money):} 두 번째로 중요한 자금 조달원입니다. 단기 담보 부채(예: 환매 조건부 시장을 통한 차입)부터 장기 무담보 부채(예: 장기 선순위 채권 발행)까지 다양합니다.
    \item \textbf{은행 자본 (Bank Equity Capital):} 은행 소유주로부터 조달된 자금입니다.
    \begin{itemize}
        \item \textbf{납입 자본 (Paid-in Capital):} 보통주 및 우선주 발행을 통해 조달됩니다.
        \item \textbf{유보 이익 (Retained Earnings):} JP Morgan 자본의 대부분은 배당되지 않고 사업에 재투자된 이익에서 나옵니다.
    \end{itemize}
\end{enumerate}
JP Morgan은 투자를 위해 많은 부채를 사용하며, 총 자금 조달에서 자본의 비중은 약 8\%입니다. 이는 은행의 자기자본 승수(Equity Multiplier)가 약 12배라는 것을 의미합니다.

\begin{warningbox}
\textbf{은행의 높은 레버리지:} 은행은 일반 기업에 비해 부채 비율이 매우 높습니다. 이는 예금이라는 저렴하고 안정적인 자금원을 활용하기 때문이지만, 동시에 금융 위기 시 취약성을 높이는 요인이 되기도 합니다.
\end{warningbox}

\subsubsection{미국 상업은행 산업 전체의 통합 대차대조표 (2018년)}
2018년 기준 미국 내 약 5,000개 상업은행의 통합 대차대조표를 살펴보겠습니다. 정부는 은행 시스템의 안전성을 모니터링하기 위해 이러한 통계를 수집합니다.

\paragraph{자산 측면:}
\begin{itemize}
    \item \textbf{총 자산:} 약 16.5조 달러.
    \item \textbf{대출 (Loans):} 약 9조 달러 (총 자산의 약 55\%).
    \begin{itemize}
        \item 상업 및 주거용 부동산 대출, 기업 대출, 개인 대출 순으로 중요합니다.
    \end{itemize}
    \item \textbf{투자 증권 (Investment Securities):} 총 자산의 약 25\%로 두 번째로 중요한 자산입니다.
    \begin{itemize}
        \item 미국 국채, 정부 기관 모기지 담보 증권 등이 가장 큰 하위 구성 요소입니다. 이 자산들은 쉽게 매각할 수 있어 은행의 현금 준비금에 대한 좋은 백업 역할을 합니다.
    \end{itemize}
    \item \textbf{현금 (Cash):} 총 자산의 약 10\%를 차지합니다.
\end{itemize}

\begin{tipbox}
\textbf{현금 보유량 증가:} 2008년 금융 위기 이후 많은 은행이 현금 부족을 겪었기 때문에, 은행들은 역사적으로 높은 수준의 현금을 보유하는 경향이 있습니다. 또한, 2008년 이후 낮은 금리는 현금 보유의 기회비용을 낮추는 요인이 되었습니다.
\end{tipbox}

\paragraph{가장 작은 지역 은행 (Community Banks)의 대출 포트폴리오:}
총 자산 10억 달러 미만의 지역 소매 은행들은 다음과 같은 특징을 보입니다.
\begin{itemize}
    \item \textbf{신용카드 역할 미미:} 신용카드는 이들 은행 사업에서 거의 역할을 하지 않습니다. 단위당 대출 발생 및 서비스 비용이 너무 비싸기 때문입니다. 대형 은행과 전문 금융 회사가 대규모 신용카드 풀을 다룹니다.
    \item \textbf{기업 대출 비중 낮음:} 기업 대출의 역할도 예상보다 작습니다. 대부분의 상업 대출은 대형 은행이 담당합니다.
    \item \textbf{부동산 대출 집중:} 지역 은행들은 주거용 및 상업용 부동산, 건설 및 개발 대출 등 부동산 대출에 집중합니다. 이들은 지역 시장을 매우 잘 알고 있으며, 이것이 이들의 비교 우위입니다.
\end{itemize}

\paragraph{부채 측면:}
\begin{itemize}
    \item \textbf{주요 자금 조달원:} 예금, 차입금, 자기 자본.
    \item \textbf{예금 (Deposits):} 총 투자의 약 77\%를 조달합니다. 대부분 미국 내 국내 지점에서 보유됩니다.
    \item \textbf{비예금 차입금 (Non-deposit Borrowings):} 자산의 약 11\%를 차지합니다.
    \item \textbf{자기 자본 (Equity Capital):} 총 자산에서 부채를 뺀 나머지 약 11\%를 조달합니다.
\end{itemize}

\begin{figure}[h!]
\centering
\caption{미국 상업은행의 부채 구성 변화 (1970년대~현재)}
\label{fig:us_bank_liabilities_trend}
% 이미지 파일이 없으므로 텍스트로 설명합니다.
\begin{tcolorbox}[title=차트 설명: 미국 상업은행의 부채 구성 변화]
이 차트는 1970년대부터 현재까지 미국 상업은행의 부채 구성을 보여줍니다.
\begin{itemize}
    \item \textbf{예금의 중요성:} 예금은 총 상업은행 자산의 70\%에서 80\% 사이를 꾸준히 유지하며 가장 중요한 자금 조달원임을 보여줍니다.
    \item \textbf{2008년 금융 위기 이후:} 2008년 금융 위기 이후 은행들은 더욱 신중해졌으며, 안정적이고 신뢰할 수 있는 예금의 역할이 더욱 중요해졌습니다.
\end{itemize}
\end{tcolorbox}
\end{figure}

\begin{summarybox}
\textbf{대차대조표 분석 요약:}
\begin{itemize}
    \item \textbf{자산 측면:} 대출, 증권, 현금이 주요 항목입니다.
    \item \textbf{부채 측면:} 예금, 자본 시장 차입금, 자기 자본이 주요 자금 조달원입니다.
    \item \textbf{최근 동향:} 2008년 금융 위기 이후 은행들은 더 많은 현금을 보유하고, 투자를 위해 더 많은 예금을 활용하는 등 더욱 신중하게 행동하고 있습니다.
\end{itemize}
\end{summarybox}

\newpage
\section{대차대조표 외 활동 (Off-Balance Sheet Activities)}

\subsection{우발 자산 및 우발 부채의 개념}
대차대조표 외 활동(Off-Balance Sheet, OBS)은 현재 은행의 대차대조표에 자산이나 부채로 기록되지 않지만, 특정 사건이 발생할 경우 은행의 재무제표나 현금 흐름에 중대한 영향을 미칠 수 있는 항목들을 말합니다. 이는 현대 회계 기준에 따라 투자자와 규제 당국에 정보를 제공하기 위해 재무제표에 기록되어야 합니다.

\begin{itemize}
    \item \textbf{우발 자산 (Contingent Asset):} 특정 사건 발생 시 은행의 자산으로 전환되거나 수입에 영향을 미치는 항목입니다.
    \begin{itemize}
        \item \textbf{예시: 신용카드 약정 (Credit Card Commitment):} 은행은 고객에게 신용 한도만큼 대출을 제공하겠다는 약속을 합니다. 고객이 신용카드를 사용할 때(우발 사건), 이 미사용 약정은 대차대조표 외 자산에서 대차대조표 내 자산(소비자 대출)으로 전환됩니다.
    \end{itemize}
    \item \textbf{우발 부채 (Contingent Liability):} 특정 사건 발생 시 은행의 부채로 전환되거나 비용에 영향을 미치는 항목입니다.
    \begin{itemize}
        \item \textbf{예시: 보증 (Guarantees) 또는 보험:} 은행이 보증을 제공하거나 보험을 판매할 때, 이는 대차대조표 외 부채로 기록됩니다. 보증 사건이 발생하거나 보험 청구가 들어오면(우발 사건), 은행은 현금 지출을 해야 하므로 대차대조표에 중대한 영향을 미칩니다.
    \end{itemize}
\end{itemize}
은행이 신용카드를 발행하거나 보험을 판매하는 것을 선호하는 이유는 이러한 활동이 선불 수수료(Upfront Fees)를 발생시키기 때문입니다.

\begin{summarybox}
\textbf{핵심 요약:} 대차대조표 외 활동은 미래의 특정 사건에 따라 은행의 자산이나 부채로 변할 수 있는 잠재적 항목입니다. 이는 은행의 현재 재무 상태를 넘어선 잠재적 위험과 수익 기회를 보여줍니다.
\end{summarybox}

\subsection{주요 대차대조표 외 활동}
은행의 대차대조표 외 활동은 전통적인 은행 대출과 유사한 활동도 있고, 보험과 유사한 활동도 있습니다.

\begin{itemize}
    \item \textbf{자산 측면의 주요 항목:}
    \begin{itemize}
        \item \textbf{미사용 대출 약정 (Undrawn Loan Commitments):} 신용카드 한도와 같이 미래에 인출될 수 있는 대출 약정입니다.
        \item \textbf{롱 파생상품 포지션 (Long Derivatives Positions):} 미래에 이익 또는 손실을 발생시킬 수 있는 파생상품 계약입니다.
    \end{itemize}
    \item \textbf{부채 측면의 주요 항목:}
    \begin{itemize}
        \item \textbf{숏 파생상품 포지션 (Short Derivatives Positions):} 미래에 이익 또는 손실을 발생시킬 수 있는 파생상품 계약입니다.
        \item \textbf{보증 (Guarantees):} 다양한 형태를 취하지만, 대부분 보험과 유사하게 작동합니다. 특정 조건 하에 은행이 보험금 지급을 해야 할 수 있습니다.
    \end{itemize}
\end{itemize}

\subsubsection{미사용 대출 약정 (Loan Commitments)}
대출 약정은 특정 최대 한도 내에서 미리 정해진 이자율로 자금을 제공하겠다는 법적 구속력이 있는 약속입니다. 이는 가계나 기업에 제공될 수 있으며, 무담보 신용카드나 주택 담보 신용 한도 등이 있습니다.

\begin{itemize}
    \item \textbf{은행의 의무:} 은행은 약정 기간 내에 언제든지 자금을 제공할 준비가 되어 있어야 합니다. 고객이 카드를 사용하면 대출이 은행의 대차대조표에 기록됩니다.
    \item \textbf{활용률 (Utilization Rates):} FDIC 자료에 따르면, 주택 담보 대출, 기업 대출, 신용카드 등 다양한 유형의 대출 약정 활용률은 100\% 미만입니다. 이는 대차대조표 외에 미인출 금액이 존재한다는 것을 의미합니다.
    \item \textbf{신용카드 예시:} 신용카드 미인출 잔액은 인출 잔액의 4~5배에 달할 수 있습니다. 은행은 이러한 미인출 잔액에 대해서도 수수료를 받을 수 있습니다.
    \item \textbf{규제 당국의 요구:} 규제 당국은 과도한 인출이 은행에 위험을 초래할 수 있으므로, 은행이 이러한 미인출 잔액을 재무제표에 공개하도록 요구합니다.
\end{itemize}

\subsubsection{파생상품 (Derivatives)}
은행은 파생상품을 사용하여 금리 변동과 같은 불리한 위험 노출을 줄일 수 있습니다. 또한, 고객과의 파생상품 거래에서 상대방이 되어 수수료를 받을 수도 있습니다.

\begin{itemize}
    \item \textbf{현금 흐름 영향:} 파생상품의 매수 및 매도는 미래에 현금 유입 또는 유출로 이어질 수 있으므로, 은행은 이러한 거래를 대차대조표 외에 기록해야 합니다.
    \item \textbf{명목 금액 (Notional Amounts):} 파생상품 보고는 명목 금액으로 이루어집니다. 금리 파생상품의 경우, 150조 달러에 달하는 엄청난 규모를 보입니다. 이는 '보험'으로 생각했을 때, 보험에 가입된 항목의 총 가치를 나타냅니다.
    \item \textbf{시장 가치와의 차이:} 명목 금액은 실제 은행의 잠재적 현금 흐름 노출이나 포지션의 시장 가치와는 다릅니다. (예: 100만 달러 주택 보험의 명목 금액은 100만 달러이지만, 보험료(시장 가치)는 2,000달러일 수 있습니다.)
    \item \textbf{성장률:} 지난 30년간 대출 약정은 은행 대차대조표와 거의 비슷한 5배 성장을 보인 반면, 파생상품은 20배 폭발적인 성장을 보였습니다.
    \item \textbf{집중도:} 이러한 파생상품 활동의 대부분은 JP Morgan Chase, Citi Bank, Bank of America, Goldman Sachs, Morgan Stanley, Wells Fargo, HSBC와 같은 소수의 대형 금융기관에 집중되어 있습니다. 이는 이 사업에서 규모의 경제가 작용함을 반영합니다.
\end{itemize}

\paragraph{JP Morgan Chase의 파생상품 노출 (2020년):}
JP Morgan은 금리 파생상품, 신용 파생상품, 외환, 주식, 상품 계약 등 수천 건의 파생상품 거래에 참여합니다. 대부분의 활동은 고객을 대신하여 이루어집니다.

\begin{itemize}
    \item \textbf{상쇄 (Netting):} JP Morgan이 금리 옵션을 동시에 매수하고 매도하는 것처럼, 많은 거래가 상쇄(Netting)될 수 있습니다. 거래가 많을수록 상쇄 효과가 커집니다.
    \item \textbf{신용부도스왑 (Credit Default Swap, CDS):} CDS는 채권이나 대출의 부도 위험에 대한 보험 계약으로 생각할 수 있습니다.
    \begin{itemize}
        \item JP Morgan은 약 1조 달러 규모의 CDS에 관여하지만, 이는 보험 판매와 보험 매수 사이에 균등하게 분산되어 있습니다.
        \item \textbf{순 노출 (Net Exposure):} 전체 순 노출은 훨씬 작습니다. JP Morgan이 매수하는 보험에 대해 지불하는 금액보다 보험 판매로 징수하는 수수료가 훨씬 적다면, 이 활동은 은행에 너무 많은 위험 노출을 만들지 않으면서 좋은 이익을 가져다줄 수 있습니다.
    \end{itemize}
\end{itemize}

\begin{warningbox}
\textbf{2008년 금융 위기와 OBS 위험:} 대차대조표 외에 숨겨진 위험은 2008년 금융 위기의 핵심 원인이었습니다. 제대로 이해되지 않은 이러한 위험은 세계 최대 금융기관들의 파산, 부실 인수 또는 구제금융으로 이어졌습니다. 이는 수많은 새로운 규제를 낳았습니다.
\end{warningbox}

\begin{summarybox}
\textbf{대차대조표 외 활동 요약:}
\begin{itemize}
    \item \textbf{개념:} 미래의 특정 사건 발생 시 은행의 재무 상태에 영향을 미칠 수 있는 잠재적 자산 및 부채입니다.
    \item \textbf{주요 활동:} 대출 약정(신용카드 한도 등)과 파생상품이 가장 중요하며, 은행에 많은 수수료 수입을 가져다줍니다.
    \item \textbf{위험:} 명목 금액이 크더라도 상쇄 효과(Netting)로 실제 순 노출은 작을 수 있지만, 2008년 금융 위기에서 보았듯이 OBS 활동의 위험은 매우 중요하며 엄격한 규제가 필요합니다.
\end{itemize}
\end{summarybox}

\newpage
\section*{FAQ (자주 묻는 질문)}

\begin{itemize}
    \item \textbf{Q: 투자은행과 상업은행의 가장 큰 차이점은 무엇인가요?}
    \item \textbf{A:} 상업은행은 주로 예금을 받아 대출을 해주는 전통적인 역할을 합니다. 반면 투자은행은 기업의 주식/채권 발행을 통한 자금 조달, 인수합병 자문, 증권 거래 등 자본시장 관련 활동에 특화되어 있습니다. 최근에는 많은 대형 은행들이 이 두 가지 기능을 모두 수행하는 경향이 있습니다.

    \item \textbf{Q: 2008년 금융 위기가 투자은행 산업에 어떤 영향을 미쳤나요?}
    \item \textbf{A:} 2008년 금융 위기는 투자은행 산업에 엄청난 변화를 가져왔습니다. Lehman Brothers는 파산했고, Bear Stearns와 Merrill Lynch는 상업은행에 인수되었습니다. Goldman Sachs와 Morgan Stanley는 생존을 위해 상업은행 지주회사로 전환해야 했습니다. 이로 인해 독립형 투자은행은 거의 사라지고, 대부분 상업은행 지주회사 산하로 편입되었습니다.

    \item \textbf{Q: '뱅킹 북'과 '트레이딩 북'은 왜 구분하나요?}
    \item \textbf{A:} 뱅킹 북은 은행의 전통적인 장기 대출 및 예금 활동을, 트레이딩 북은 단기적인 증권 매매 활동을 기록합니다. 이 둘을 구분하는 이유는 활동의 성격(장기 vs. 단기), 위험 프로필, 수익 인식 방식(이자 수익 vs. 거래 손익)이 다르기 때문입니다. 이는 은행의 다양한 사업 모델을 명확히 보여줍니다.

    \item \textbf{Q: 대차대조표 외 활동(OBS)이 중요한 이유는 무엇인가요?}
    \item \textbf{A:} OBS 활동은 현재 은행의 대차대조표에는 나타나지 않지만, 미래에 특정 사건이 발생하면 은행의 자산, 부채, 현금 흐름에 중대한 영향을 미칠 수 있는 잠재적 위험과 수익 기회를 포함합니다. 신용카드 약정이나 파생상품 등이 대표적이며, 2008년 금융 위기에서 OBS 위험이 은행 시스템에 얼마나 큰 영향을 미칠 수 있는지 드러나면서 그 중요성이 더욱 강조되었습니다.

    \item \textbf{Q: 파생상품의 '명목 금액'이 크다고 해서 은행의 위험이 그만큼 큰 것은 아닌가요?}
    \textbf{A:} 명목 금액은 파생상품 계약의 기초자산 총액을 나타내지만, 실제 은행의 잠재적 손실 노출과는 다릅니다. 은행은 수많은 파생상품 거래를 통해 롱 포지션과 숏 포지션을 동시에 보유하여 위험을 상쇄(Netting)할 수 있습니다. 따라서 순 노출(Net Exposure)은 명목 금액보다 훨씬 작을 수 있습니다. 그러나 명목 금액의 급격한 증가는 잠재적 위험의 증가를 시사하므로 규제 당국이 면밀히 주시합니다.
\end{itemize}

\newpage
\section*{빠르게 훑어보기 (1페이지 요약)}

\begin{tcolorbox}[title=\textbf{투자은행의 역할과 조직}]
\begin{itemize}
    \item \textbf{역할:} 기업의 자금 조달(주식/채권 발행), M\&A 자문, 증권 거래를 통해 차입자와 저축자를 연결.
    \item \textbf{주요 사업:} 자금 조달(인수), 증권 거래(시장 조성/자기 계정 거래), M\&A 자문.
    \item \textbf{조직:} JP Morgan 사례처럼 상업은행 지주회사 내 '뱅킹'과 '마켓'으로 구성.
    \item \textbf{유형:} 상업은행 계열사(JP Morgan), 전문화된 투자은행(Goldman Sachs), 독립형(Lazard).
    \item \textbf{산업 변화:} 1990년대 규제 완화와 2008년 금융 위기로 대부분 상업은행과 통합.
\end{itemize}
\end{tcolorbox}

\begin{tcolorbox}[title=\textbf{투자은행 활동의 세부 사항}]
\begin{itemize}
    \item \textbf{증권 인수:} IPO, 추가 발행, 채권 인수, 신디케이트론 주간사 역할. 북빌딩을 통해 가격 결정 및 판매.
    \item \textbf{M\&A 자문:} 기업 가치 평가, 협상 지원, 자금 조달 지원. 높은 수수료 수익원.
    \item \textbf{증권 거래:}
    \begin{itemize}
        \item \textbf{시장 조성:} 매수/매도 호가 제시로 유동성 공급, 비드-애스크 스프레드 수익.
        \item \textbf{자기 계정 거래:} 은행 자체 이익을 위한 포지션 거래 및 차익 거래.
    \end{itemize}
    \item \textbf{수익 비중:} 거래 수익(약 75\%)이 투자은행 수수료(약 25\%)보다 훨씬 큼.
\end{itemize}
\end{tcolorbox}

\begin{tcolorbox}[title=\textbf{현대 은행의 재무제표 분석}]
\begin{itemize}
    \item \textbf{대차대조표:} 자산(은행의 투자)과 부채(자금 조달원)를 보여주는 재무 스냅샷.
    \item \textbf{뱅킹 북 vs. 트레이딩 북:}
    \begin{itemize}
        \item \textbf{뱅킹 북:} 전통적 장기 대출/예금 활동.
        \item \textbf{트레이딩 북:} 단기 증권/상품/통화 거래.
    \end{itemize}
    \item \textbf{JP Morgan 자산:} 대출(가장 큼), 투자 증권, 현금, 연방 기금 판매.
    \item \textbf{JP Morgan 부채:} 예금(가장 큼), 차입금, 자기 자본. 높은 레버리지(자본 8\%, 승수 12배).
    \item \textbf{산업 전체:} 대출(55\%), 투자 증권(25\%), 현금(10\%). 예금(77\%)이 주요 자금원.
    \item \textbf{지역 은행:} 부동산 대출에 집중, 신용카드/기업 대출 비중 낮음.
\end{itemize}
\end{tcolorbox}

\begin{tcolorbox}[title=\textbf{대차대조표 외 활동 (OBS)}]
\begin{itemize}
    \item \textbf{개념:} 현재 대차대조표에 없지만, 특정 사건 발생 시 재무에 영향을 미치는 우발 자산/부채.
    \item \textbf{주요 활동:}
    \begin{itemize}
        \item \textbf{미사용 대출 약정:} 신용카드 한도 등. 인출 시 대차대조표 내 자산으로 전환.
        \item \textbf{파생상품:} 위험 헤지 또는 거래 목적. 명목 금액은 크지만 순 노출은 상쇄 효과로 작을 수 있음.
    \end{itemize}
    \item \textbf{중요성:} 선불 수수료 발생. 2008년 금융 위기의 핵심 원인 중 하나로, 잠재적 위험 관리가 중요.
\end{itemize}
\end{tcolorbox}

\newpage

\newpage
\section*{개요}

\begin{summarybox}[title={핵심 요약}]
이 문서는 은행의 손익계산서(Income Statement)를 분석하여 현대 은행이 어떻게 수익을 창출하는지 설명합니다. 전통적인 대출 사업에서 발생하는 이자수익과 이자비용의 차이인 순이자이익(Net Interest Income)뿐만 아니라, 투자은행 업무, 자산 관리 등 다양한 비즈니스 라인에서 발생하는 비이자수익(Non-Interest Revenue)의 중요성을 강조합니다. JP Morgan Chase 사례 연구를 통해 이러한 수익원들이 어떻게 결합되어 은행의 최종 순이익을 구성하는지 구체적으로 살펴봅니다.
\end{summarybox}

\newpage
\section*{용어 정리}

\begin{table}[h!]
\centering
\caption{주요 용어 설명}
\label{tab:glossary}
\begin{adjustbox}{width=\textwidth,center}
\begin{tabular}{lllc}
\toprule
\textbf{용어 (한글)} & \textbf{쉬운 설명} & \textbf{원어} & \textbf{비고} \\
\midrule
스프레드 & 대출로 벌어들인 이자수익과 예금에 지급한 이자비용의 차이. & Spread & 은행의 전통적인 수익원 \\
순이자이익 & 이자수익에서 이자비용을 뺀 금액. 은행 대출 사업의 핵심 수익성 지표. & Net Interest Income (NII) & 손익계산서의 최상단 지표 \\
비이자수익 & 이자 외의 모든 수익. 수수료, 거래 이익 등이 포함됨. & Non-Interest Revenue & 현대 은행의 중요한 수익원 \\
비이자비용 & 이자 지급과 무관한 모든 운영 비용. 급여, 임대료, 기술 투자 등. & Non-Interest Expense & 은행 운영에 필요한 제반 비용 \\
자기매매 & 은행이 고객을 대신하는 것이 아니라, 자체 자금으로 금융 시장에서 & Proprietary Trading & 고위험 고수익 전략 \\
& 투기적 포지션을 취하여 이익을 얻는 활동. & & \\
인수 업무 & 기업의 주식이나 채권 발행을 돕고 이를 투자자에게 판매하는 업무. & Underwriting & 투자은행의 핵심 기능 \\
시장 조성 & 특정 증권의 매수-매도 호가를 제시하여 시장 유동성을 공급하는 역할. & Market Making & 거래 활성화에 기여 \\
3-6-3 모델 & 과거 은행업의 단순한 수익 모델. 예금에 3\% 지급, 6\%로 대출, 3\% 이익. & 3-6-3 Model & 1980년대 이후 압박받음 \\
\bottomrule
\end{tabular}
\end{adjustbox}
\end{table}

\newpage
\section*{핵심 개념·원리: 은행의 손익계산서}

은행의 손익계산서는 단순히 이자수익과 이자비용의 차이만을 보여주는 것이 아닙니다. 현대 은행은 복잡하고 다각화된 사업 구조를 가지고 있으며, 다양한 방식으로 수익을 창출하고 비용을 지출합니다.

\begin{tcolorbox}[mybox,title={은행의 기본적인 수익 구조: 스프레드}]
\textbf{1. 한 줄 핵심 요약:} 은행은 예금을 받아 대출해주고, 대출 이자수익과 예금 이자비용의 차이인 '스프레드'를 통해 이익을 얻습니다.

\textbf{2. 직관적 예시/비유:} 친구에게 돈을 빌려주고(예금 수취) 다른 친구에게 더 높은 이자로 빌려주는(대출) 중개인과 같습니다. 빌려준 돈에서 받은 이자(대출 이자)와 빌린 돈에 대해 지불한 이자(예금 이자)의 차이가 당신의 이익이 됩니다.

\textbf{3. 기술적 정확한 설명:} 상업은행은 예금을 수취하여 개인, 기업, 부동산 대출에 활용합니다. 이 과정에서 은행은 자금을 차입하고, 대출하며, 이자수익(interest income)과 이자비용(interest expense)의 차이인 스프레드(spread)를 얻습니다.
\end{tcolorbox}

하지만 현대 은행의 수익 구조는 이보다 훨씬 복잡합니다. 오늘날의 은행은 대규모의 다각화된 사업체이며, 대출 사업에서 발생하는 이자수익 외에도 다양한 사업 라인에서 수수료 수익(fee revenue)을 창출합니다.

\subsection*{손익계산서의 세 가지 주요 구성 요소}

은행의 손익계산서는 크게 세 부분으로 나눌 수 있습니다. 이 세 가지 구성 요소를 이해하면 현대 은행이 어떻게 이익을 창출하는지 명확하게 파악할 수 있습니다.

\begin{enumerate}
    \item \textbf{대출 사업의 수익성 (순이자이익):}
    \begin{tcolorbox}[mybox,title={순이자이익 (Net Interest Income, NII)}]
    \textbf{1. 한 줄 핵심 요약:} 은행의 전통적인 대출 및 투자 활동에서 발생하는 이자수익에서 자금 조달에 드는 이자비용을 뺀 금액입니다.

    \textbf{2. 직관적 예시/비유:} 당신이 부동산 임대 사업을 한다고 가정해 봅시다. 임대료 수입(이자수익)에서 대출 이자(이자비용)를 뺀 금액이 순수하게 부동산에서 벌어들인 돈과 같습니다.

    \textbf{3. 기술적 정확한 설명:}
    \begin{itemize}
        \item \textbf{이자수익 (Interest Income):} 은행이 대출을 제공하거나 유가증권에 투자함으로써 얻는 수익입니다. 여기에는 개인 대출, 기업 대출, 부동산 대출, 다른 금융기관에 대한 대출, 그리고 은행이 보유한 유가증권(예: 국채, 회사채) 투자에서 발생하는 이자가 포함됩니다.
        \item \textbf{이자비용 (Interest Expense):} 은행이 예금자에게 지급하거나 금융시장에서 자금을 차입할 때 발생하는 비용입니다. 여기에는 고객 예금에 대한 이자, 단기 및 장기 차입금에 대한 이자, 그리고 은행이 발행한 장기 채권의 쿠폰 지급액 등이 포함됩니다.
    \end{itemize}
    \textbf{순이자이익 (NII) = 이자수익 - 이자비용} 이며, 이는 은행의 대출 사업 수익성을 나타내는 가장 중요한 지표입니다.
    \end{tcolorbox}

    \item \textbf{비이자수익 (Non-Interest Revenue):}
    \begin{tcolorbox}[mybox,title={비이자수익 (Non-Interest Revenue)}]
    \textbf{1. 한 줄 핵심 요약:} 이자수익 외에 은행이 다양한 서비스 제공을 통해 얻는 모든 수수료 및 기타 수익입니다.

    \textbf{2. 직관적 예시/비유:} 식당을 운영한다고 할 때, 음식 판매 수익(이자수익) 외에 배달 수수료, 예약 수수료, 케이터링 서비스 수수료 등 부가적인 서비스로 버는 돈과 같습니다.

    \textbf{3. 기술적 정확한 설명:} 현대 은행은 대출 외에도 다양한 금융 서비스를 제공하며 수수료를 받습니다. 이는 '수수료 수익(fee income)'이라고도 불리며, 투자은행 업무, 자산 관리, 중개 서비스, 계좌 관리 수수료, 신용카드 수수료 등 매우 광범위합니다. 이는 은행의 수익원을 다각화하고 안정성을 높이는 데 기여합니다.
    \end{tcolorbox}

    \item \textbf{비이자비용 (Non-Interest Expense):}
    \begin{tcolorbox}[mybox,title={비이자비용 (Non-Interest Expense)}]
    \textbf{1. 한 줄 핵심 요약:} 이자 지급과 직접적인 관련이 없는 은행 운영에 필요한 모든 비용을 포괄합니다.

    \textbf{2. 직관적 예시/비유:} 회사를 운영할 때, 원자재 구매 비용(이자비용) 외에 직원 급여, 사무실 임대료, 전기세, 마케팅 비용 등 회사를 유지하는 데 드는 모든 일반적인 지출과 같습니다.

    \textbf{3. 기술적 정확한 설명:} 비이자비용은 이자 지급과 무관한 모든 운영 비용을 포함하는 포괄적인 항목입니다. 여기에는 직원 급여 및 보너스(가장 큰 비중), 사무실 임대료, 기술 투자 비용, 마케팅 비용, 감가상각비, 그리고 기타 소규모 운영 비용 등이 포함됩니다.
    \end{tcolorbox}
\end{enumerate}

이 세 가지 요소를 종합하여 은행의 최종 \textbf{순이익 (Net Income)}이 계산됩니다.
\textbf{순이익 = 순이자이익 + 비이자수익 - 비이자비용}

\newpage
\section*{사례 연구: JP Morgan Chase의 손익 구조 (2020년 기준)}

JP Morgan Chase는 세계에서 가장 크고 다각화된 은행 중 하나입니다. 이 은행의 2020년 손익계산서를 통해 위에서 설명한 개념들이 실제 어떻게 적용되는지 구체적으로 살펴보겠습니다.

\subsection*{JP Morgan의 사업 부문}

JP Morgan의 사업은 크게 두 가지로 나눌 수 있습니다.

\begin{enumerate}
    \item \textbf{전통적인 은행 업무 (대출):}
    소매 금융(retail banking) 부문과 도매 금융(wholesale business) 내의 상업 대출(commercial lending) 부문이 여기에 해당합니다. 이 부문에서 주로 \textbf{순이자이익}이 발생합니다.

    \item \textbf{비이자수익 창출 업무 (수수료):}
    주로 투자은행(investment banking) 부문에서 발생하며, 인수(underwriting) 업무, 시장 조성(market making), 자산 관리(asset management) 등이 포함됩니다. 여기서 \textbf{비이자수익}이 발생합니다.
\end{enumerate}

\subsection*{이자수익과 이자비용}

2020년 JP Morgan의 이자 관련 수익과 비용은 다음과 같습니다.

\begin{itemize}
    \item \textbf{이자수익 (Interest Income):} 약 \textbf{650억 달러}
    \begin{itemize}
        \item 대출 포트폴리오 (loan portfolio)
        \item 은행 계정(banking book) 및 트레이딩 포트폴리오(trading portfolio)에 보유된 투자 유가증권
        \item 다른 금융기관에 대한 대출
    \end{itemize}
    \item \textbf{이자비용 (Interest Expense):}
    \begin{itemize}
        \item 예금에 대한 이자
        \item 기타 차입금 (단기 및 장기 부채, JP Morgan이 발행한 장기 채권의 쿠폰 지급액 포함)
    \end{itemize}
\end{itemize}

이 둘의 차이인 \textbf{순이자이익 (Net Interest Income)}은 2020년에 약 \textbf{550억 달러}였습니다. 이는 JP Morgan의 대출 사업에서 발생한 수익성을 나타냅니다.

\subsection*{비이자수익 (Non-Interest Income)}

순이자이익에서 전체 수익성으로 넘어가기 위해서는 수수료 수익을 고려해야 합니다. JP Morgan의 비이자수익은 매우 다양하며, 주로 도매 금융 부문에서 발생합니다. 2020년 JP Morgan의 비이자수익은 약 \textbf{640억 달러}에 달했습니다. 주요 항목은 다음과 같습니다.

\begin{enumerate}
    \item \textbf{투자은행 수수료 (Investment Banking Fees):} 약 \textbf{90억 달러}
    \begin{itemize}
        \item \textbf{M\&A 자문 (M\&A advisory):} 기업 인수합병 관련 자문 서비스.
        \item \textbf{주식 인수 (equity underwriting):} 기업의 주식 발행을 돕고 이를 투자자에게 판매.
        \item \textbf{채권 및 신디케이트론 인수 (bond and syndicated loan underwriting):} 기업의 채권 발행 및 여러 은행이 공동으로 제공하는 대출(신디케이트론) 주선.
    \end{itemize}
    \item \textbf{자기매매 이익 (Proprietary Trading / Principal Transactions):} 약 \textbf{180억 달러}
    \begin{itemize}
        \item JP Morgan의 트레이더들이 고객을 대신하는 것이 아니라, \textbf{은행 자체 계정(on its own account)}으로 금융 시장에서 투기적 포지션을 취하여 이익을 창출하는 활동입니다. 이는 고위험 고수익 사업으로 분류될 수 있습니다.
    \end{itemize}
    \item \textbf{자산 관리 및 중개 서비스 수수료 (Asset Management and Brokerage Services Fees):} 약 \textbf{180억 달러}
    \begin{itemize}
        \item JP Morgan이 고객의 금융 시장 투자를 돕는 서비스입니다.
        \item \textbf{자산 관리 수수료 (asset management fees):} 고객 포트폴리오를 관리해주는 대가로 받는 수수료.
        \item \textbf{중개 수수료 (brokerage commission fees):} 고객의 거래를 중개해주고 받는 수수료.
    \end{itemize}
    \item \textbf{기타 수수료:} 예금 계좌 및 모기지 관련 수수료 (예: 초과 인출 수수료, 연체료), 신용카드 수수료 등.
\end{enumerate}

\begin{tcolorbox}[warningbox,title={주목할 점: 비이자수익의 비중}]
JP Morgan과 같이 크고 다각화된 은행의 경우, 2020년 비이자수익(약 640억 달러)이 대출 사업에서 발생한 순이자이익(약 540억 달러)보다 실제로 더 컸습니다. 이는 현대 은행의 수익 구조에서 비이자수익이 얼마나 중요한 부분을 차지하는지 보여줍니다.
\end{tcolorbox}

\subsection*{비이자비용 (Non-Interest Expense)}

JP Morgan의 비이자비용은 다음과 같은 항목들을 포함합니다.

\begin{itemize}
    \item \textbf{급여 및 보너스 (Salaries and bonuses):} 가장 큰 비중을 차지합니다.
    \item 임대료 (Rent)
    \item 기술 투자 (Technology)
    \item 마케팅 (Marketing)
    \item 기타 소규모 항목들
\end{itemize}
은행은 또한 세금(taxes)을 납부해야 합니다.

\subsection*{최종 순이익 (Net Income)}

JP Morgan의 2020년 순이익은 약 \textbf{300억 달러}였습니다. 이는 순이자이익과 비이자수익을 합산한 후 비이자비용을 차감하여 계산된 은행의 최종 수익입니다.

\newpage
\section*{현대 은행업의 수익 구조 변화}

JP Morgan의 사례는 현대 은행이 단순히 대출과 예금의 이자 차이로만 수익을 내는 것이 아님을 명확히 보여줍니다.

\subsection*{수수료 수익의 중요성}

\begin{tcolorbox}[mybox,title={수수료 수익 증가의 배경}]
\textbf{1. 한 줄 핵심 요약:} 21세기 은행들은 낮은 금리 환경과 대출 시장 경쟁 심화로 인해 이자수익 의존도를 줄이고 수수료 수익을 통해 수익원을 다각화하고 있습니다.

\textbf{2. 직관적 예시/비유:} 스마트폰 제조사가 단순히 기기 판매(이자수익) 외에 앱 스토어 수수료, 클라우드 서비스 구독료(비이자수익) 등으로 수익을 창출하는 것과 유사합니다. 핵심 제품의 수익성이 낮아지면 부가 서비스로 보완하는 것입니다.

\textbf{3. 기술적 정확한 설명:}
\begin{itemize}
    \item \textbf{낮은 금리 환경:} 지속적으로 낮은 금리는 은행의 순이자마진(Net Interest Margin)을 압박하여 대출 사업의 수익성을 저하시킵니다.
    \item \textbf{대출 시장 경쟁:} 핀테크 기업 등 새로운 경쟁자들의 등장으로 대출 시장의 경쟁이 심화되어 전통적인 대출 수익만으로는 성장이 어려워졌습니다.
    \item \textbf{다각화의 필요성:} 수수료 수익은 이러한 환경에서 은행이 잃어버린 이자수익을 보충하고, 수익원을 다각화하여 사업 안정성을 높이는 중요한 수단이 됩니다.
\end{itemize}
\end{tcolorbox}

\subsection*{3-6-3 모델과 수수료 비즈니스로의 전환}

1980년대에 '3-6-3 모델'이라는 전통적인 은행업 모델이 큰 압박을 받기 시작하면서, 은행들은 수수료 기반 비즈니스로의 전환을 가속화했습니다. 3-6-3 모델은 예금에 3\% 이자를 지급하고, 6\% 이자로 대출을 해주며, 3\%의 이익을 남기는 단순한 구조를 의미했습니다. 하지만 규제 완화와 경쟁 심화로 이러한 단순한 모델로는 더 이상 수익을 내기 어려워졌습니다.

\subsection*{미국 은행의 평균 수익 구성}

JP Morgan은 미국 최대 은행으로서 다각화된 사업 구조를 가지고 있지만, 수수료 수익에 의존하는 경향은 비단 JP Morgan만의 이야기가 아닙니다. 미국 은행 전체를 평균적으로 보면, 은행 총수익의 약 3분의 1이 수수료에서 발생합니다. 이는 수수료 수익이 현대 은행업에서 보편적으로 중요한 역할을 하고 있음을 보여줍니다.

\newpage
\section*{결론 및 요약}

\begin{summarybox}[title={핵심 정리}]
이 세션에서는 은행의 손익계산서를 심층적으로 분석했습니다. 우리는 현대 은행이 전통적인 대출 사업에서 발생하는 \textbf{순이자이익}뿐만 아니라, 다양한 금융 서비스에서 발생하는 \textbf{비이자수익}을 통해 이익을 창출한다는 것을 확인했습니다. 특히, JP Morgan Chase 사례를 통해 비이자수익이 순이자이익보다 더 큰 비중을 차지할 수 있음을 보았습니다. 21세기 은행들은 낮은 금리와 치열한 경쟁 환경 속에서 수익원을 다각화하고 안정성을 확보하기 위해 수수료 기반 비즈니스의 중요성이 더욱 커지고 있습니다.
\end{summarybox}

\end{document}
